% =====================================================
% part3_vectorspace.tex ― 第3部:ベクトル空間と線形写像
% =====================================================
\part{ベクトル空間と線形写像}

% =====================================================
\chapter{ベクトル空間 ― 「線形」の舞台を抽象化する}\label{chap:vs}
% =====================================================

\begin{goal}
\begin{itemize}
\item ベクトル空間の $8$ つの公理を言葉で理解し、$\R^n$ 以外の集合
――行列・多項式・関数・数列――が、同じ公理のもとでベクトル空間に
なることを確認できるようになる。
\item 公理\textbf{だけ}を使って、零ベクトルの一意性・$0\vect{v} = \vect{0}$・
$(-1)\vect{v} = -\vect{v}$ といった「当たり前に見える」事実を導けるようになる。
\item ある集合と演算が\textbf{ベクトル空間にならない}ことを、
どの公理が破れるかを指し示して説明できるようになる。
\end{itemize}
\end{goal}

\section{動機 ― なぜ「舞台」を抽象化するのか}

第II部の最後で、連立一次方程式の解の構造定理が
微分方程式や漸化式にもそっくり現れることを見た。
「そっくり」で済ませず、\textbf{文字どおり同じ定理}として一度で証明したい。
そのために必要なのが、舞台そのものの抽象化――ベクトル空間の公理である。

そもそも第I部・第II部で $\R^n$ を扱ったとき、私たちが本当に使っていた性質は
何だったろうか。線形結合、線形独立、解空間――どれを振り返っても、
登場するのは\textbf{和とスカラー倍、そしてそれらの計算規則}だけである。
「成分が縦に $n$ 個並んでいること」そのものは、ほとんど使っていない。
だとすれば、成分という見た目を捨て、\textbf{和とスカラー倍が同じ規則で
計算できる集合}を一括して扱えば、$\R^n$ で示した定理はすべて、その集合でも
そのまま成り立つはずである。これが抽象化のねらいである。

抽象化は難しくするための操作ではなく、\textbf{一度の証明を使い回すための節約}である。
本章で公理を据え、次章以降で基底・次元・次元定理を公理だけから証明すれば、
それらは行列の空間でも、多項式の空間でも、微分方程式の解の空間でも、
\textbf{改めて確認することなく}使える。第II部の深掘りで「似ている」と
予告した現象が、ここで「同じもの」に格上げされる。

\section{公理による定義}

$\R^n$ で実際に使った性質を振り返ると、和とスカラー倍、
およびそれらの計算規則だけであった。そこで、
\textbf{「成分」を忘れ、演算の規則だけを要求する}ことにする。

次の定義には公理が $8$ つ並ぶが、身構える必要はない。要するに、
\begin{enumerate}
\item 足し算が、順序や組み方によらずできる(結合律・交換律)、
\item 基準となる零ベクトル $\vect{0}$ と、各ベクトルの「マイナス」がある、
\item スカラー倍が、足し算とうまく噛み合う(分配律など)、
\item $1$ 倍しても変わらない、
\end{enumerate}
という四つのことを、漏れなく細かく書き下しただけである。
この四点を頭に置けば、以下の公理は\textbf{暗記すべきリストではなく、
「普通の計算」を抽象化したもの}に見えてくる。

\begin{teigi}[ベクトル空間]
集合 $V$ に、和 $\vect{u} + \vect{v}$ とスカラー倍 $c\vect{u}$
($c \in \R$)が定義されていて、次の公理を満たすとき、
$V$ を($\R$ 上の)\textbf{ベクトル空間}といい、$V$ の要素を\textbf{ベクトル}と呼ぶ。

任意の $\vect{u}, \vect{v}, \vect{w} \in V$ と $a, b \in \R$ に対して:
\begin{enumerate}
\item $(\vect{u} + \vect{v}) + \vect{w} = \vect{u} + (\vect{v} + \vect{w})$
\item $\vect{u} + \vect{v} = \vect{v} + \vect{u}$
\item 零ベクトル $\vect{0}$ が存在し、$\vect{u} + \vect{0} = \vect{u}$
\item 各 $\vect{u}$ に逆ベクトル $-\vect{u}$ が存在し、$\vect{u} + (-\vect{u}) = \vect{0}$
\item $a(\vect{u} + \vect{v}) = a\vect{u} + a\vect{v}$
\item $(a + b)\vect{u} = a\vect{u} + b\vect{u}$
\item $(ab)\vect{u} = a(b\vect{u})$
\item $1\vect{u} = \vect{u}$
\end{enumerate}
\end{teigi}

公理は多く見えるが、内容は「足し算とスカラー倍が、
$\R^n$ と同じ感覚で計算できること」の言語化にすぎない。
重要なのは、\textbf{ベクトルが「数の組」である必要はもうない}という点である。

\subsection*{公理から導かれる基本事実}

公理に書かれていない性質も、公理\textbf{だけ}から論理的に導ける。
「成分計算」を一切使わず、$8$ つの公理だけを根拠に証明する練習をしておこう。
ここで証明する事実は、以後どのベクトル空間でも無条件に使ってよい。

\begin{teiri}[公理から導かれる基本事実]\label{thm:vs-basic}
任意のベクトル空間 $V$ において、次が成り立つ。
\begin{enumerate}
\item 零ベクトルはただ一つである。
\item 任意の $\vect{v} \in V$ に対し、スカラー $0$ 倍は零ベクトル:$0\vect{v} = \vect{0}$。
\item 任意の $\vect{v} \in V$ に対し、$(-1)\vect{v}$ は $\vect{v}$ の逆ベクトル:$(-1)\vect{v} = -\vect{v}$。
\end{enumerate}
\end{teiri}

\noindent\textbf{(証明)}\quad
(1) $\vect{0}$ と $\vect{0}'$ がともに零ベクトルの公理3を満たすとする。
$\vect{0}'$ が零ベクトルだから $\vect{0} = \vect{0} + \vect{0}'$、
さらに公理2(交換律)と $\vect{0}$ が零ベクトルであることから
$\vect{0} + \vect{0}' = \vect{0}' + \vect{0} = \vect{0}'$。
両者を合わせて $\vect{0} = \vect{0}'$。零ベクトルは一意である。

(2) 公理6で $a = b = 0$ とすると
\[
0\vect{v} = (0 + 0)\vect{v} = 0\vect{v} + 0\vect{v}.
\]
両辺に $0\vect{v}$ の逆ベクトル $-(0\vect{v})$ を加える。左辺は公理4より $\vect{0}$、
右辺は公理1(結合律)と公理4・公理3より
$0\vect{v} + \bigl(0\vect{v} + (-(0\vect{v}))\bigr) = 0\vect{v} + \vect{0} = 0\vect{v}$。
したがって $\vect{0} = 0\vect{v}$。

(3) 公理8・公理6と(2)より
\[
\vect{v} + (-1)\vect{v} = 1\vect{v} + (-1)\vect{v} = \bigl(1 + (-1)\bigr)\vect{v}
= 0\vect{v} = \vect{0}.
\]
すなわち $(-1)\vect{v}$ は $\vect{v}$ に足すと零になる。よって
\[
(-1)\vect{v} = \vect{0} + (-1)\vect{v} = (-\vect{v} + \vect{v}) + (-1)\vect{v}
= -\vect{v} + \bigl(\vect{v} + (-1)\vect{v}\bigr) = -\vect{v} + \vect{0} = -\vect{v}. \qquad\blacksquare
\]

\begin{chui}
(2)では「スカラーの $0$」と「ベクトルの $\vect{0}$」という\textbf{別物}が、
$0\vect{v} = \vect{0}$ という形で結びついた。記号が似ているので混同しやすいが、
左辺の $0$ は数、右辺の $\vect{0}$ はベクトルである。この区別を意識すると、
抽象的な議論の見通しがよくなる。
\end{chui}

\section{例 ― 思いがけないものがベクトルになる}

どの例でも、公理の確認は最終的に\textbf{実数の計算規則}に帰着する。
各例につき一つだけ、公理を実際に確かめてみせる。

\begin{rei}[数ベクトル空間と行列の空間]
$\R^n$ は成分ごとの演算でベクトル空間である(これが原型)。
零ベクトルは $\vect{0} = (0, \dots, 0)$、$\vect{u} = (u_1, \dots, u_n)$ の
逆ベクトルは $(-u_1, \dots, -u_n)$ である。
$m \times n$ 行列全体の集合 $M_{m,n}$ も、行列の和とスカラー倍でベクトル空間になる。\\
\textbf{公理の確認(公理5).}\ $M_{m,n}$ で $c(A + B) = cA + cB$ は、
$(i,j)$ 成分を比べると $c(a_{ij} + b_{ij}) = c\,a_{ij} + c\,b_{ij}$、
すなわち\textbf{実数の分配法則}そのものである。他の公理も同様に成分ごとの
実数の規則に帰着する。零ベクトルは零行列 $O$ である。
\end{rei}

\begin{rei}[多項式の空間]
$n$ 次以下の実係数多項式全体
\[
\R[x]_n = \{\, a_0 + a_1 x + \cdots + a_n x^n \mid a_0, \dots, a_n \in \R \,\}
\]
は、多項式の和とスカラー倍でベクトル空間になる。
たとえば $(1 + 2x) + (3 - x) = 4 + x$、$3(1 + 2x) = 3 + 6x$。
次数の制限を外した多項式全体 $\R[x]$ もベクトル空間である。\\
\textbf{公理の確認(公理6).}\ $p = 1 + 2x$ で $(a + b)p = ap + bp$ を確かめる。
左辺は $(a + b)(1 + 2x) = (a + b) + 2(a + b)x$、
右辺は $a(1 + 2x) + b(1 + 2x) = (a + b) + (2a + 2b)x$ で一致する。
零ベクトルは零多項式(すべての係数が $0$)である。
\end{rei}

\begin{rei}[関数の空間]
実数全体で定義された連続関数全体の集合 $C(\R)$ は、
$(f + g)(x) = f(x) + g(x)$、$(cf)(x) = c\,f(x)$ という自然な演算で
ベクトル空間になる。$\sin x$ も $e^x$ も「ベクトル」である。\\
\textbf{公理の確認(公理2).}\ 交換律 $f + g = g + f$ は、各点 $x$ で
$(f + g)(x) = f(x) + g(x) = g(x) + f(x) = (g + f)(x)$、
すなわち\textbf{実数の和の交換律}に帰着する。零ベクトルは恒等的に $0$ をとる関数、
$f$ の逆ベクトルは $-f$(各点で符号を変えた関数)である。
\end{rei}

\begin{rei}[数列の空間]
実数列 $(a_0, a_1, a_2, \dots)$ 全体も、項ごとの和とスカラー倍で
ベクトル空間になる。「無限次元版の $\R^n$」といえる。\\
\textbf{公理の確認(公理3・公理4).}\ 零ベクトルはすべての項が $0$ の数列
$(0, 0, 0, \dots)$ で、任意の $(a_n)$ にこれを足しても変わらない。
$(a_n)$ の逆ベクトルは $(-a_n)$ であり、和は零数列になる。
項ごとに実数の世界の話をしているだけである。
\end{rei}

\begin{itembox}[l]{深掘り:なぜわざわざ公理化するのか}
抽象化の利得は\textbf{「一度証明すれば、すべての例で成り立つ」}ことに尽きる。
これから証明する定理(基底・次元・次元定理など)は、
公理だけから導かれるので、$\R^n$ でも多項式でも関数でも数列でも、
\textbf{確認なしに}適用できる。第II部の深掘りで予告した
「微分方程式の解の構造が連立方程式と同じ」は、もはや類似ではなく、
\textbf{同じ舞台(関数のベクトル空間)上の同じ定理}になる。

もう一つの利得は、議論に使ってよい性質が公理に限定されることで、
\textbf{証明の本質が見える}ことである。成分計算に頼った証明は
「なぜ成り立つのか」を隠すことがあるが、公理からの証明は隠せない。
公理的方法は20世紀数学の標準語であり、線形代数はその最良の入門である。
\end{itembox}

\begin{mondai}\label{q:fib}
漸化式 $a_{n+2} = a_{n+1} + a_n$ を満たす実数列全体の集合 $V$ は、
数列の空間の演算(項ごとの和・スカラー倍)に関してベクトル空間になることを示せ。
すなわち、$(a_n), (b_n) \in V$ ならば $(a_n + b_n) \in V$、
$c(a_n) \in V$ となることを確かめよ。
(フィボナッチ数列はこの $V$ の一員である。$V$ の「次元」が何になるかは、
本章を読み進めてから問 \ref{q:basis2} で考える。)
\end{mondai}

\section{ベクトル空間にならない例}

「公理を満たす」とは、$8$ つすべてを満たすことである。
たった一つでも破れれば、その集合と演算はベクトル空間ではない。
\textbf{破れる公理を一つ指し示す}のが、否定の最短証明である。

\begin{hanrei}[演算が閉じていない]\label{ex:halfplane}
上半平面 $H = \{\,(x, y) \in \R^2 \mid y \ge 0\,\}$ に、$\R^2$ の標準の
和とスカラー倍を入れる。和については閉じている($y$ 成分が非負どうしの和は非負)が、
\textbf{スカラー倍で閉じていない}。実際 $(0, 1) \in H$ に対し
\[
(-1)(0, 1) = (0, -1) \notin H.
\]
演算の結果が集合の外に出てしまうので、そもそも「$H$ 上の演算」になっておらず、
$H$ はベクトル空間ではない。
\end{hanrei}

\begin{hanrei}[和の定義を変えると公理が壊れる]\label{ex:weirdadd}
台集合は $\R^2$ のまま、和だけを
\[
(x_1, y_1) \oplus (x_2, y_2) = (x_1 + x_2,\ 0)
\]
と定義し直す(スカラー倍は標準のまま)。一見「足し算らしい」が、
\textbf{零ベクトルが存在しない}。もし零ベクトル $\vect{z} = (a, b)$ があれば、
任意の $(x, y)$ に対して $(x, y) \oplus \vect{z} = (x, y)$ のはずだが、
左辺は $(x + a,\ 0)$ で第 $2$ 成分が必ず $0$ になり、$y \ne 0$ のときは
$(x, y)$ に等しくならない。たとえば $(1, 2) \oplus (a, b) = (1 + a, 0) \ne (1, 2)$。
公理3が破れるので、この演算ではベクトル空間にならない。
\end{hanrei}

\begin{chui}
反例 \ref{ex:weirdadd} の教訓は、\textbf{集合が同じでも演算が変われば
ベクトル空間かどうかは変わる}ということである。「ベクトル空間」とは
集合そのものの性質ではなく、\textbf{集合と二つの演算の組}の性質である。
\end{chui}

\section{例題}

公理の使い方を、計算・検証・証明の三つの型で練習する。
以下は\textbf{解答つきの例題}である。手を動かして追ってほしい。

\begin{rei}[例題・基本(計算)]
$\R[x]_2$ で $2(1 + x - x^2) - 3(2 - x)$ を計算せよ。\\
\textbf{解.}\ 項ごとに展開して
\[
2(1 + x - x^2) = 2 + 2x - 2x^2, \qquad -3(2 - x) = -6 + 3x.
\]
和をとると $(2 - 6) + (2 + 3)x - 2x^2 = -4 + 5x - 2x^2$。
多項式の演算も、結局は係数(=実数)の計算である。
\end{rei}

\begin{rei}[例題・基本(検証)]
連続関数の空間 $C(\R)$ について、零ベクトルは何か。
また $f(x) = x^2$ の逆ベクトルを書け。\\
\textbf{解.}\ 零ベクトルは、任意の $g$ に足しても $g$ を変えない元、
すなわち\textbf{恒等的に $0$ をとる関数} $o(x) = 0$ である。
$f$ の逆ベクトルは $f$ に足すと $o$ になる元だから、各点で符号を変えた
$(-f)(x) = -x^2$ である。「定数 $0$」という数ではなく「$0$ という関数」が
零ベクトルである点に注意。
\end{rei}

\begin{rei}[例題・基本(証明)]\label{rei:cancel}
任意のベクトル空間で\textbf{消去律}
「$\vect{u} + \vect{v} = \vect{u} + \vect{w}$ ならば $\vect{v} = \vect{w}$」
が成り立つことを、公理だけから示せ。\\
\textbf{解.}\ 仮定の両辺に左から $-\vect{u}$ を加える。公理1(結合律)で組み替えると
\[
(-\vect{u} + \vect{u}) + \vect{v} = (-\vect{u} + \vect{u}) + \vect{w}.
\]
公理4より $-\vect{u} + \vect{u} = \vect{0}$、公理3より $\vect{0} + \vect{v} = \vect{v}$、
$\vect{0} + \vect{w} = \vect{w}$ だから $\vect{v} = \vect{w}$。
成分を一切使っていないことに注目せよ。
\end{rei}

\begin{rei}[例題・標準(検証)]\label{rei:exotic}
正の実数全体 $\R_{>0}$ に、和とスカラー倍を
\[
x \oplus y = xy, \qquad c \odot x = x^c
\]
で定める($\oplus$ は通常の掛け算、$\odot$ は累乗)。これがベクトル空間に
なることを確かめ、零ベクトルと、$x$ の逆ベクトルを求めよ。\\
\textbf{解.}\ 零ベクトル $e$ は $x \oplus e = xe = x$ より $e = 1$。
$x$ の逆ベクトル $x'$ は $x \oplus x' = x x' = 1$ より $x' = 1/x$
($x > 0$ だから存在する)。主な公理を確かめると
\[
c \odot (x \oplus y) = (xy)^c = x^c y^c = (c \odot x) \oplus (c \odot y),
\]
\[
(a + b) \odot x = x^{a+b} = x^a x^b = (a \odot x) \oplus (b \odot x), \qquad
1 \odot x = x^1 = x.
\]
他の公理も同様に成り立つ。\textbf{見た目は掛け算と累乗なのに、
公理のうえでは立派なベクトル空間}である。
\end{rei}

\begin{chui}
例題 \ref{rei:exotic} は、いわば\textbf{腕試しの例}である。
これを見て「ベクトル空間は何でもありなのか」と思うのは早とちりで、
事情はむしろ逆である。$\oplus, \odot$ をどう定めようとも、
\textbf{$8$ つの公理をすべて満たさなければ}ベクトル空間とは呼べない
(満たさない例は次節の反例で見る)。この例が面白いのは、
一見ベクトルらしくない正の実数が、うまい演算のもとで\textbf{公理を全部満たしてしまう}
点にある。まずは標準的な例($\R^n$・多項式・関数)で公理に慣れ、
この種の例は「公理さえ満たせばよい」と納得したあとに味わうとよい。
\end{chui}

\begin{rei}[例題・標準(計算と関係)]
$C(\R)$ の三つの元 $f(x) = \cos^2 x$、$g(x) = \sin^2 x$、$h(x) = 1$
(定数関数)を考える。$f + g$ を $h$ で表せ。\\
\textbf{解.}\ 各点 $x$ で $(f + g)(x) = \cos^2 x + \sin^2 x = 1 = h(x)$ だから、
$f + g = h$、すなわち $f + g - h = o$(零ベクトル)。
見た目の異なる三つの関数のあいだに線形の関係 $f + g - h = o$ が隠れている。
このような関係の有無が、次章の「線形独立・従属」の主題になる。
\end{rei}

\begin{rei}[例題・標準(証明)]\label{rei:czero}
任意のベクトル空間で、任意のスカラー $c$ に対し $c\vect{0} = \vect{0}$ が
成り立つことを公理だけから示せ。\\
\textbf{解.}\ 公理3より $\vect{0} + \vect{0} = \vect{0}$ だから、公理5を使って
\[
c\vect{0} = c(\vect{0} + \vect{0}) = c\vect{0} + c\vect{0}.
\]
両辺に $c\vect{0}$ の逆ベクトルを加えれば、定理 \ref{thm:vs-basic}(2)と同じ要領で
$\vect{0} = c\vect{0}$ を得る。
これで「数の $0$ 倍」も「ベクトル $\vect{0}$ の何倍」も、ともに $\vect{0}$ に
なることが分かった。
\end{rei}

\begin{matome}
\begin{itemize}
\item \textbf{ベクトル空間}とは、和とスカラー倍が $8$ つの公理を満たす
「集合と演算の組」である。ベクトルは数の組とは限らない。
\item 行列・多項式・関数・数列は、いずれも自然な演算でベクトル空間になる。
公理の確認は最終的に\textbf{実数の計算規則}に帰着する。
\item 零ベクトルの一意性・$0\vect{v} = \vect{0}$・$(-1)\vect{v} = -\vect{v}$・
$c\vect{0} = \vect{0}$・消去律は、すべて公理\textbf{だけ}から導かれる
(成分計算は不要)。
\item ベクトル空間に\textbf{ならない}ことは、破れる公理を一つ示せばよい
(閉じていない/零ベクトルがない、など)。
\item 「ベクトル空間か」は集合だけでなく\textbf{演算の定め方}にも依存する。
\end{itemize}
\end{matome}

\section*{章末問題}

\begin{mondai}[基本]\label{q:vs-ex1}
$\R[x]_2$ で $3(2 - x + x^2) - 2(1 + 2x)$ を計算せよ。
\end{mondai}

\begin{mondai}[基本]\label{q:vs-ex2}
連続関数の空間 $C(\R)$ の零ベクトルは何か。
また $f(x) = e^x$ の逆ベクトル $-f$ を具体的に書け。
\end{mondai}

\begin{mondai}[基本]\label{q:vs-ex3}
$\R^2$ に標準の和とスカラー倍を入れる。次の部分集合は、その演算で
閉じている(ベクトル空間になる)か。閉じていないものには理由を一つ述べよ。
\[
L_1 = \{\,(x, y) \mid y = 2x \,\}, \qquad
L_2 = \{\,(x, y) \mid y = 2x + 1 \,\}.
\]
\end{mondai}

\begin{mondai}[基本]\label{q:vs-ex4}
$2 \times 2$ 実行列の空間で、$A = \mat{1 & 2 \\ 0 & -1}$,
$B = \mat{3 & 0 \\ 1 & 2}$ に対し $2A - B$ を計算せよ。
また、この空間の零ベクトルは何か。
\end{mondai}

\begin{mondai}[標準]\label{q:vs-ex5}
$\R^2$ に、和は標準のまま、スカラー倍だけを
$c \odot (x, y) = (cx, 0)$ と定義し直す。$8$ つの公理のうち破れるものを
一つ挙げ、具体的なベクトルで破れることを示せ。
\end{mondai}

\begin{mondai}[標準]\label{q:vs-ex6}
漸化式 $a_{n+2} = 2a_{n+1} - a_n$ を満たす実数列全体 $W$ が、
数列空間の演算(項ごとの和・スカラー倍)に関してベクトル空間になることを示せ。
\end{mondai}

\begin{mondai}[標準]\label{q:vs-ex7}
$C(\R)$ の中で、偶関数全体
$E = \{\, f \mid f(-x) = f(x)\ (\forall x) \,\}$ が、関数空間の演算で
ベクトル空間になることを示せ。
\end{mondai}

\begin{mondai}[発展]\label{q:vs-ex8}
例題 \ref{rei:exotic} のベクトル空間 $\R_{>0}$
($x \oplus y = xy$, $c \odot x = x^c$)で、$\vect{u} = 2$, $\vect{v} = 8$ に対し
$(3 \odot \vect{u}) \oplus \bigl((-1) \odot \vect{v}\bigr)$ を計算せよ。
結果はこの空間の何にあたるか。
\end{mondai}

\begin{mondai}[発展]\label{q:vs-ex9}
$2 \times 2$ 実行列の空間で、対称行列全体
$S = \{\, A \mid {}^t\!A = A \,\}$ と交代行列全体
$T = \{\, A \mid {}^t\!A = -A \,\}$ が、それぞれ標準の演算でベクトル空間に
なることを示せ。さらに $S \cap T$ を求めよ。
\end{mondai}

\begin{mondai}[証明]\label{q:vs-ex10}
任意のベクトル空間で、$\vect{v} + \vect{v} = \vect{v}$ を満たすベクトルは
$\vect{v} = \vect{0}$ に限ることを、公理だけから示せ。
\end{mondai}

\begin{mondai}[証明]\label{q:vs-ex11}
スカラー $a \ne 0$ とベクトル $\vect{v}$ について、$a\vect{v} = \vect{0}$ ならば
$\vect{v} = \vect{0}$ であることを公理だけから示せ
($c\vect{0} = \vect{0}$ は用いてよい)。
\end{mondai}

\bigskip
{\small
\noindent\textbf{【章末問題の方針】}\quad
まず自力で取り組み、行き詰まったら下の方針を見よ。記述による完全解答は
巻末「問の解答」にある。
\begin{itemize}
\item \textbf{問 \ref{q:vs-ex1}}:各項をスカラー倍してから、同じ次数の係数をそろえて足す。
\item \textbf{問 \ref{q:vs-ex2}}:零ベクトルは「足しても相手を変えない関数」、
逆ベクトルは「足すと零関数になる関数」。数ではなく関数で答える。
\item \textbf{問 \ref{q:vs-ex3}}:原点を含むか、和・スカラー倍で集合の外に出ないかを調べる。
切片のある直線に注意。
\item \textbf{問 \ref{q:vs-ex4}}:成分ごとに計算する。零ベクトルは零行列。
\item \textbf{問 \ref{q:vs-ex5}}:スカラー倍を変えたので公理5〜8を疑う。
$1 \odot \vect{v}$ がもとに戻るか確かめる。
\item \textbf{問 \ref{q:vs-ex6}}:和とスカラー倍が同じ漸化式を保つ(=$W$ に留まる)
ことだけ示せばよい。
\item \textbf{問 \ref{q:vs-ex7}}:偶関数の定義 $f(-x) = f(x)$ に戻り、
$f + g$ と $cf$ が偶関数になることを確かめる。
\item \textbf{問 \ref{q:vs-ex8}}:$\odot$ は累乗、$\oplus$ は掛け算。定義どおり計算し、
結果を零ベクトル($= 1$)と見比べる。
\item \textbf{問 \ref{q:vs-ex9}}:転置の線形性を使う。共通部分は
${}^t\!A = A$ と ${}^t\!A = -A$ を連立する。
\item \textbf{問 \ref{q:vs-ex10}}:$\vect{v} = \vect{v} + \vect{0}$ と書いて消去律。
ベクトルでは割り算をしない。
\item \textbf{問 \ref{q:vs-ex11}}:$a \ne 0$ なので $\dfrac{1}{a}$ を掛ける。
$\vect{v} = 1\vect{v} = \left(\dfrac{1}{a}\cdot a\right)\vect{v}$ から出発する。
\end{itemize}
}

% =====================================================
\chapter{部分空間と生成}\label{chap:subspace}
% =====================================================

\begin{goal}
\begin{itemize}
\item 部分空間の定義(和とスカラー倍で閉じた部分集合)を述べ、
与えられた集合が部分空間かどうかを、根拠を挙げて判定・証明できるようになる。
\item 部分空間から新しい部分空間を作る二つの方法――\textbf{共通部分}と
\textbf{和空間}――を使えるようになり、和集合ではだめな理由を
反例で説明できるようになる。
\item ベクトルの組が\textbf{生成する(張る)}部分空間を理解し、
「このベクトルは張られた空間に入るか」「この組は空間全体を張るか」を
掃き出し法・行列式の計算で判定できるようになる。
\end{itemize}
\end{goal}

\section{動機 ― 大きな空間の中の「閉じた世界」}

前章で、行列・多項式・関数・数列と、思いがけない集合たちが
ベクトル空間になることを見た。ここで前章の問を振り返ると、
共通のパターンに気づく。偶関数全体(問 \ref{q:vs-ex7})は
関数空間 $C(\R)$ の一部、漸化式を満たす数列全体
(問 \ref{q:vs-ex6}・問 \ref{q:fib})は数列空間の一部、
対称行列全体(問 \ref{q:vs-ex9})は行列空間の一部だった。
そしてどの場合も、$8$ つの公理を一から確かめてはいない。
実際に確かめたのは、\textbf{和とスカラー倍でその集合の外に出ないこと}
――たったこれだけだった。

同じ現象は第II部にもあった。同次方程式 $A\vect{x} = \vect{0}$ の解全体は、
「解の和も解、解のスカラー倍も解」(定理 \ref{thm:homog-closed})という
意味で閉じた $\R^n$ の部分集合である。

本章はこのパターンに\textbf{部分空間}という名前を与え、
「なぜ閉じていることの確認だけで、公理が全部ついてくるのか」を
はっきりさせる。さらに、部分空間どうしを組み合わせる方法
(共通部分・和空間)と、指定したベクトルたちから部分空間を作る方法
(生成)を整備する。次章の主役「基底と次元」が住んでいるのは、
ほとんどの場合、この部分空間たちの世界である。

\section{部分空間 ― 定義と判定法}

\begin{teigi}[部分空間]
ベクトル空間 $V$ の空でない部分集合 $W$ が、
\begin{enumerate}
\item $\vect{u}, \vect{v} \in W \implies \vect{u} + \vect{v} \in W$
\item $\vect{u} \in W,\ c \in \R \implies c\vect{u} \in W$
\end{enumerate}
を満たすとき、$W$ を $V$ の\textbf{部分空間}という。
\end{teigi}

条件 1・2 を合わせて「$W$ は和とスカラー倍について\textbf{閉じている}」
という。定義には $8$ つの公理が一つも現れないが、次の定理のとおり、
閉じてさえいれば公理はすべて「親」の $V$ から無料でついてくる。

\begin{teiri}[部分空間はそれ自身ベクトル空間である]\label{thm:sub-inherit}
$W$ を $V$ の部分空間とすると、次が成り立つ。
\begin{enumerate}
\item $\vect{0} \in W$。
\item $\vect{u} \in W$ ならば $-\vect{u} \in W$。
\item $W$ は、$V$ の和とスカラー倍(を $W$ に制限したもの)によって、
それ自身ベクトル空間になる。
\end{enumerate}
\end{teiri}

\noindent\textbf{(証明)}\quad
(1) $W$ は空でないから、あるベクトル $\vect{u} \in W$ が取れる。
条件2を $c = 0$ で使うと $0\vect{u} \in W$、そして
定理 \ref{thm:vs-basic}(2) より $0\vect{u} = \vect{0}$。
よって $\vect{0} \in W$ である。

(2) 条件2を $c = -1$ で使うと $(-1)\vect{u} \in W$、
定理 \ref{thm:vs-basic}(3) より $(-1)\vect{u} = -\vect{u}$ である。

(3) まず条件 1・2 により、和とスカラー倍は結果が $W$ からはみ出さない、
$W$ の中で完結した演算になっている。公理のうち
1, 2, 5, 6, 7, 8 は「$V$ の\textbf{すべての}ベクトルについて成り立つ等式」
だから、$W$ のベクトルに限っても当然成り立つ($V$ から遺伝する)。
残るは「存在」を主張する公理 3・4 だが、公理3の零ベクトルは (1) で、
公理4の逆ベクトルは (2) で、どちらも $W$ の中に確保できた。 $\blacksquare$

\begin{chui}[判定の実際]\label{chui:sub-check}
与えられた集合 $W$ が部分空間かどうかは、次の順に調べるのが速い。
\begin{enumerate}
\item まず $\vect{0} \in W$ か。含まなければその時点で部分空間ではない
(定理 \ref{thm:sub-inherit}(1) の対偶)。
\textbf{原点を含まない集合は部分空間になりえない}。
\item $\vect{0} \in W$ なら、和とスカラー倍の閉性を\textbf{文字を使って}
確かめる(具体的な数値で一例を確かめただけでは証明にならない)。
\end{enumerate}
ただし $\vect{0} \in W$ は\textbf{必要条件にすぎない}。
原点を含んでいても閉じていない集合はある(例題 \ref{rei:sub-curve})。
\end{chui}

\begin{rei}
\begin{itemize}
\item $\R^3$ の中で、原点を通る直線と原点を通る平面は部分空間である。
原点を通らない直線・平面は部分空間でない。
\item 同次方程式 $A\vect{x} = \vect{0}$ の解空間は $\R^n$ の部分空間である
(第II部の議論そのもの)。
\item $\R[x]_2$ は $\R[x]_5$ の部分空間である。
\item $C(\R)$ の中で、微分方程式 $y'' + y = 0$ の解全体は部分空間である
(解の和も解、解の定数倍も解)。
\end{itemize}
\end{rei}

\begin{rei}[自明な部分空間]
どんなベクトル空間 $V$ にも、零ベクトルだけの集合 $\{\vect{0}\}$ と
$V$ 自身という二つの部分空間が必ずある。実際
$\vect{0} + \vect{0} = \vect{0}$、$c\vect{0} = \vect{0}$
(例題 \ref{rei:czero})だから $\{\vect{0}\}$ は閉じており、
$V$ 自身が閉じているのは定義そのものである。
$\{\vect{0}\}$ は最小の、$V$ は最大の部分空間であり、
他のすべての部分空間はこの両極のあいだにある。
\end{rei}

\begin{mondai}\label{q:subsp}
次の $\R^3$ の部分集合のうち、部分空間であるものをすべて選び、
そうでないものには理由を一つ述べよ。
\[
W_1 = \{\, (x, y, z) \mid x + y + z = 0 \,\}, \quad
W_2 = \{\, (x, y, z) \mid x + y + z = 1 \,\}, \quad
W_3 = \{\, (x, y, z) \mid x \ge 0 \,\}
\]
\end{mondai}

\section{反例 ― 破れ方はひととおりではない}

\begin{hanrei}[和集合は部分空間とは限らない]\label{hanrei:sub-union}
$\R^2$ の $x$ 軸 $L_1 = \{(x, 0) \mid x \in \R\}$ と
$y$ 軸 $L_2 = \{(0, y) \mid y \in \R\}$ は、どちらも部分空間である。
しかし和集合 $L_1 \cup L_2$(両軸を合わせた「十字」)は部分空間ではない。
スカラー倍では閉じている(軸上の点を何倍しても同じ軸の上)のに、
\textbf{和で閉じていない}:
\[
(1, 0) + (0, 1) = (1, 1) \notin L_1 \cup L_2.
\]
部分空間どうしを単純に「合併」しても部分空間にはならないのである。
合併の正しい代役は、本章で後述する\textbf{和空間}である。
\end{hanrei}

\begin{hanrei}[和で閉じてもスカラー倍で破れる]\label{hanrei:sub-lattice}
座標がともに整数である点の全体(格子)
\[
\mathbb{Z}^2 = \{\, (m, n) \mid m, n \text{ は整数} \,\}
\]
は、原点を含み、和についても閉じている(整数どうしの和は整数)。
しかしスカラー倍では閉じていない:
\[
\frac{1}{2}\,(1, 0) = \left(\frac{1}{2},\ 0\right) \notin \mathbb{Z}^2.
\]
前章の反例 \ref{ex:halfplane}(上半平面)は $-1$ 倍で破れ、
反例 \ref{hanrei:sub-union} は和で破れ、この例は $\frac{1}{2}$ 倍で破れる。
破れる場所は集合ごとに違うから、条件 1・2 は\textbf{両方とも}
確認しなければならない。
\end{hanrei}

\section{共通部分と和空間}

部分空間が二つあるとき、そこから新しい部分空間を作る操作を二つ用意する。

\begin{teiri}[共通部分は部分空間]\label{thm:sub-cap}
$W_1, W_2$ が $V$ の部分空間ならば、共通部分 $W_1 \cap W_2$ も
$V$ の部分空間である。
\end{teiri}

\noindent\textbf{(証明)}\quad
$\vect{0}$ は $W_1$ にも $W_2$ にも属するから、$W_1 \cap W_2$ は空でない。
$\vect{u}, \vect{v} \in W_1 \cap W_2$ とすると、
$\vect{u}, \vect{v} \in W_1$ で $W_1$ は閉じているから
$\vect{u} + \vect{v} \in W_1$。同様に $\vect{u} + \vect{v} \in W_2$。
よって $\vect{u} + \vect{v} \in W_1 \cap W_2$ である。
スカラー倍についてもまったく同様に、$c\vect{u}$ は $W_1$ と $W_2$ の
両方に属する。 $\blacksquare$

\begin{rei}[二つの平面の共通部分]\label{rei:sub-cap}
$\R^3$ の二つの部分空間(原点を通る平面)
\[
W_1 = \{\, (x, y, z) \mid x + y + z = 0 \,\}, \qquad
W_2 = \{\, (x, y, z) \mid x - y + z = 0 \,\}
\]
の共通部分を求めよう。両方の式を連立し、辺々引くと $2y = 0$、
すなわち $y = 0$。このとき残る条件は $x + z = 0$ だから
\[
W_1 \cap W_2 = \{\, t\,(1,\ 0,\ -1) \mid t \in \R \,\}
\]
――原点を通る直線である(検算:$(1, 0, -1)$ は
$1 + 0 - 1 = 0$、$1 - 0 - 1 = 0$ を確かに満たす)。
「共通部分を求める」ことは「条件式を連立して解く」ことにほかならず、
計算は第II部の掃き出し法の世界に帰着する。
\end{rei}

反例 \ref{hanrei:sub-union} で見たとおり、和集合 $W_1 \cup W_2$ は
部分空間になるとは限らない。「$W_1$ と $W_2$ の両方を含む部分空間」が
ほしければ、和集合の代わりに次を使う。

\begin{teigi}[和空間]\label{def:sub-sum}
$V$ の部分空間 $W_1, W_2$ に対して
\[
W_1 + W_2 = \{\, \vect{w}_1 + \vect{w}_2 \mid
\vect{w}_1 \in W_1,\ \vect{w}_2 \in W_2 \,\}
\]
を $W_1$ と $W_2$ の\textbf{和空間}という。
\end{teigi}

\begin{teiri}[和空間は部分空間]\label{thm:sub-sum}
$W_1 + W_2$ は $V$ の部分空間であり、しかも
$W_1 \cup W_2$ を含む部分空間の中で\textbf{最小}のものである
(すなわち、$W_1 \cup W_2$ を含むどんな部分空間 $W$ も
$W_1 + W_2$ を丸ごと含む)。
\end{teiri}

\noindent\textbf{(証明)}\quad
閉性:$\vect{w}_1 + \vect{w}_2$ と $\vect{w}'_1 + \vect{w}'_2$
($\vect{w}_i, \vect{w}'_i \in W_i$)の和は、公理1・2で並べ替えて
\[
(\vect{w}_1 + \vect{w}'_1) + (\vect{w}_2 + \vect{w}'_2)
\]
と書け、各 $W_i$ の閉性より第 $1$ 項は $W_1$ に、第 $2$ 項は $W_2$ に
属するから、これは再び $W_1 + W_2$ の元である。スカラー倍も
$c(\vect{w}_1 + \vect{w}_2) = c\vect{w}_1 + c\vect{w}_2 \in W_1 + W_2$。
また $\vect{0} = \vect{0} + \vect{0} \in W_1 + W_2$ で空でない。
なお $\vect{w}_1 = \vect{w}_1 + \vect{0}$ と書けば
$W_1 \subseteq W_1 + W_2$ が分かり、$W_2$ についても同様である。

最小性:$W$ を $W_1 \cup W_2$ を含む部分空間とする。任意の
$\vect{w}_1 + \vect{w}_2 \in W_1 + W_2$ について、
$\vect{w}_1, \vect{w}_2 \in W$ であり、$W$ は和で閉じているから
$\vect{w}_1 + \vect{w}_2 \in W$。よって $W_1 + W_2 \subseteq W$ である。
$\blacksquare$

\begin{rei}
$\R^2$ で $L_1 = x$ 軸、$L_2 = y$ 軸とすると、
$(x, y) = (x, 0) + (0, y)$ よりすべての点が両軸の元の和で書けるから
$L_1 + L_2 = \R^2$ である。反例 \ref{hanrei:sub-union} の「十字」
$L_1 \cup L_2$ を含む最小の部分空間は、平面全体だったのである。
\end{rei}

\section{生成 ― 指定したベクトルから最小の部分空間を作る}

部分空間のもう一つの作り方は、「このベクトルたちを含む空間がほしい」
という要望から出発するものである。含むべきベクトルを決めたら、
閉性を保つために必要なもの(それらの線形結合)をすべて追加する
――それが生成である。

\begin{teigi}[生成される部分空間]
$\vect{a}_1, \dots, \vect{a}_k \in V$ の線形結合全体の集合
\[
\langle \vect{a}_1, \dots, \vect{a}_k \rangle
= \{\, c_1\vect{a}_1 + \cdots + c_k\vect{a}_k \mid c_1, \dots, c_k \in \R \,\}
\]
は $V$ の部分空間になる。これを $\vect{a}_1, \dots, \vect{a}_k$ が
\textbf{生成する}(または\textbf{張る})部分空間という。
\end{teigi}

\begin{teiri}[生成された部分空間の最小性]\label{thm:span-min}
$\langle \vect{a}_1, \dots, \vect{a}_k \rangle$ は
$\vect{a}_1, \dots, \vect{a}_k$ をすべて含む $V$ の部分空間であり、
しかもそのような部分空間の中で\textbf{最小}のものである。
\end{teiri}

\noindent\textbf{(証明)}\quad
線形結合どうしの和は
\[
(c_1\vect{a}_1 + \cdots + c_k\vect{a}_k)
+ (d_1\vect{a}_1 + \cdots + d_k\vect{a}_k)
= (c_1 + d_1)\vect{a}_1 + \cdots + (c_k + d_k)\vect{a}_k
\]
と再び線形結合になり(公理1・2・6による並べ替え)、スカラー倍も
$c(c_1\vect{a}_1 + \cdots + c_k\vect{a}_k)
= (cc_1)\vect{a}_1 + \cdots + (cc_k)\vect{a}_k$ と線形結合になる。
係数をすべて $0$ にすれば $\vect{0}$ も入っているから、部分空間である。
各 $\vect{a}_i$ は「第 $i$ 係数だけ $1$、他は $0$」の線形結合だから
$\langle \vect{a}_1, \dots, \vect{a}_k \rangle$ に属する。

最小性:$W$ を $\vect{a}_1, \dots, \vect{a}_k$ をすべて含む部分空間とする。
$W$ はスカラー倍で閉じるから $c_i\vect{a}_i \in W$、和で閉じるから
それらを順に足した $c_1\vect{a}_1 + \cdots + c_k\vect{a}_k$ も $W$ に属する。
よって $\langle \vect{a}_1, \dots, \vect{a}_k \rangle \subseteq W$ である。
$\blacksquare$

幾何的には、$1$ 本のベクトルが張るのは直線、
独立な $2$ 本が張るのは平面である。
「$\vect{b}$ が $\vect{a}_1, \dots, \vect{a}_k$ の線形結合で書けるか」という
第I部以来の問いは、「$\vect{b} \in \langle \vect{a}_1, \dots, \vect{a}_k \rangle$ か」
と言い換えられる。

\subsection*{生成の問いは掃き出し法に帰着する}

$V = \R^n$ のときは、この言い換えのおかげで、生成に関する問いが
すべて第I・II部の計算に落ちる。
$\vect{a}_1, \dots, \vect{a}_k$ を列に並べた $n \times k$ 行列を
$A = (\vect{a}_1 \ \cdots \ \vect{a}_k)$ とすると、
$c_1\vect{a}_1 + \cdots + c_k\vect{a}_k = A\vect{c}$ だから:
\begin{itemize}
\item $\vect{b} \in \langle \vect{a}_1, \dots, \vect{a}_k \rangle$
$\iff$ 連立方程式 $A\vect{c} = \vect{b}$ が解をもつ
$\iff$ $\mathrm{rank}\,A = \mathrm{rank}\,(A \mid \vect{b})$
(定理 \ref{thm:solvable})。
\item $\langle \vect{a}_1, \dots, \vect{a}_k \rangle = \R^n$
$\iff$ どんな $\vect{b}$ に対しても $A\vect{c} = \vect{b}$ が解をもつ
$\iff$ $\mathrm{rank}\,A = n$。特に $k = n$ のときは
$\det A \neq 0$ と同値である(定理 \ref{thm:reg-rank})。
\end{itemize}
抽象的に見えた「生成」の問いが、掃き出し法と行列式という
手計算に帰着する。例題で練習しよう。

\begin{itembox}[l]{深掘り:解の集合ふたたび ― 部分空間と平行移動}
第II部の解の構造定理(定理 \ref{thm:structure})を、部分空間の言葉で
読み直してみよう。同次方程式 $A\vect{x} = \vect{0}$ の解全体 $W$ は
$\R^n$ の部分空間である。一方、非同次方程式 $A\vect{x} = \vect{b}$
($\vect{b} \neq \vect{0}$)の解全体は、一つの解 $\vect{x}_0$ を使って
\[
\vect{x}_0 + W = \{\, \vect{x}_0 + \vect{w} \mid \vect{w} \in W \,\}
\]
と書ける。これは\textbf{部分空間 $W$ を $\vect{x}_0$ だけ平行移動した集合}
である。原点を含まないから部分空間ではないが、でたらめな集合でもない
――形は $W$ のまま、位置だけずれている。このような集合を
\textbf{アフィン部分空間}と呼ぶ。問 \ref{q:vs-ex3} の直線
$L_2$($y = 2x + 1$)も、部分空間 $y = 2x$ を $(0, 1)$ だけずらした
アフィン部分空間である。

「線形の世界では、まず原点を通る部分空間を理解し、原点を通らないものは
\textbf{部分空間+平行移動}として処理する」――この二段構えは、
微分方程式の一般解(斉次解+特殊解)から第VII部の最小二乗法まで、
繰り返し現れる線形代数の基本戦略である。
\end{itembox}

\section{例題}

部分空間の判定(肯定・否定)と、生成に関する計算を練習する。
以下は\textbf{解答つきの例題}である。手を動かして追ってほしい。

\begin{rei}[例題・基本(部分空間であることの証明)]\label{rei:sub-plane}
$W = \{\, (x, y, z) \in \R^3 \mid 2x - y + 3z = 0 \,\}$ が
$\R^3$ の部分空間であることを示せ。\\
\textbf{解.}\ まず $2 \cdot 0 - 0 + 3 \cdot 0 = 0$ より $\vect{0} \in W$
(特に $W$ は空でない)。
$\vect{u} = (x_1, y_1, z_1)$, $\vect{v} = (x_2, y_2, z_2) \in W$ とすると
\[
2(x_1 + x_2) - (y_1 + y_2) + 3(z_1 + z_2)
= (2x_1 - y_1 + 3z_1) + (2x_2 - y_2 + 3z_2) = 0 + 0 = 0
\]
より $\vect{u} + \vect{v} \in W$。スカラー倍も
\[
2(cx_1) - (cy_1) + 3(cz_1) = c(2x_1 - y_1 + 3z_1) = c \cdot 0 = 0
\]
より $c\vect{u} \in W$。よって $W$ は部分空間である
(幾何的には原点を通る平面。同次方程式
$(2\ {-1}\ 3)\,\vect{x} = 0$ の解空間と見ることもできる)。
\end{rei}

\begin{rei}[例題・基本(部分空間でないことの証明)]\label{rei:sub-curve}
次の $\R^2$ の部分集合が部分空間で\textbf{ない}ことを示せ。
\[
A_1 = \{\, (x, y) \mid x + y = 2 \,\}, \qquad
A_2 = \{\, (x, y) \mid y = x^2 \,\}
\]
\textbf{解.}\ $A_1$:原点は $0 + 0 = 0 \neq 2$ より $A_1$ に属さない。
原点を含まない集合は部分空間ではない(注意 \ref{chui:sub-check})。

$A_2$:こちらは原点を含む($0 = 0^2$)。しかし
$(1, 1) \in A_2$ に対し $2\,(1, 1) = (2, 2)$ は
$2^2 = 4 \neq 2$ より $A_2$ の外に出る。スカラー倍で閉じないから
部分空間ではない(和でも破れる:$(1,1) + (-1,1) = (0, 2) \notin A_2$)。
\textbf{原点を含むことは必要条件にすぎない}ことを示す例である。
\end{rei}

\begin{rei}[例題・基本(生成された空間に入るか)]\label{rei:sub-member}
$\vect{a}_1 = \mat{1 \\ 1 \\ 0}$, $\vect{a}_2 = \mat{0 \\ 1 \\ 2}$ について、
$\vect{b} = \mat{1 \\ 3 \\ 4}$ と $\vect{b}' = \mat{3 \\ 1 \\ 4}$ が
$\langle \vect{a}_1, \vect{a}_2 \rangle$ に属するか判定せよ。\\
\textbf{解.}\ $c_1\vect{a}_1 + c_2\vect{a}_2 = \vect{b}$ を成分ごとに書くと
\[
c_1 = 1, \qquad c_1 + c_2 = 3, \qquad 2c_2 = 4.
\]
第 $1$ 式から $c_1 = 1$、第 $2$ 式から $c_2 = 2$。これは第 $3$ 式
$2 \cdot 2 = 4$ も満たすから解が存在し、
$\vect{b} = \vect{a}_1 + 2\vect{a}_2 \in \langle \vect{a}_1, \vect{a}_2 \rangle$。

$\vect{b}'$ については $c_1 = 3$, $3 + c_2 = 1$ より $c_2 = -2$ だが、
第 $3$ 式は $2c_2 = -4 \neq 4$ となり矛盾。解が存在しないから
$\vect{b}' \notin \langle \vect{a}_1, \vect{a}_2 \rangle$ である
(拡大係数行列で見れば
$\mathrm{rank}\,A = 2$, $\mathrm{rank}\,(A \mid \vect{b}') = 3$ という
定理 \ref{thm:solvable} の不成立にあたる)。
\end{rei}

\begin{rei}[例題・標準(行列の空間の部分空間)]\label{rei:sub-trace}
$2 \times 2$ 行列のうち、対角成分の和(\textbf{トレース})が $0$ の
もの全体
\[
W = \{\, A \in M_{2,2} \mid \mathrm{tr}\,A = 0 \,\}
\]
が $M_{2,2}$ の部分空間であることを示せ。\\
\textbf{解.}\ 成分で書くと、$A = (a_{ij})$, $B = (b_{ij})$ に対し
\[
\mathrm{tr}(A + B) = (a_{11} + b_{11}) + (a_{22} + b_{22})
= \mathrm{tr}\,A + \mathrm{tr}\,B, \qquad
\mathrm{tr}(cA) = c\,\mathrm{tr}\,A
\]
が成り立つ(トレースの線形性)。よって
$\mathrm{tr}\,A = \mathrm{tr}\,B = 0$ ならば
$\mathrm{tr}(A + B) = 0$、$\mathrm{tr}(cA) = 0$ であり、
零行列も $\mathrm{tr}\,O = 0$ を満たすから、$W$ は閉じており部分空間である。
なお $W$ の元は $\mat{a & b \\ c & -a}$ の形の行列全体である。
\end{rei}

\begin{rei}[例題・標準(全体を張るかの判定)]\label{rei:sub-span}
$\vect{a}_1 = \mat{1 \\ 1 \\ 1}$, $\vect{a}_2 = \mat{1 \\ 2 \\ 4}$,
$\vect{a}_3 = \mat{1 \\ 3 \\ 9}$ は $\R^3$ を生成するか。\\
\textbf{解.}\ 列に並べた行列の行列式を第 $1$ 行で余因子展開すると
\[
\det \mat{1 & 1 & 1 \\ 1 & 2 & 3 \\ 1 & 4 & 9}
= 1\,(2 \cdot 9 - 3 \cdot 4) - 1\,(1 \cdot 9 - 3 \cdot 1)
+ 1\,(1 \cdot 4 - 2 \cdot 1) = 6 - 6 + 2 = 2 \neq 0.
\]
よって行列は正則で $\mathrm{rank} = 3$。どんな $\vect{b} \in \R^3$ に
対しても $A\vect{c} = \vect{b}$ が(ただ一通りに)解けるから、
$\langle \vect{a}_1, \vect{a}_2, \vect{a}_3 \rangle = \R^3$ である。
なお $\vect{a}_j$ は $(1,\ t,\ t^2)$ に $t = 1, 2, 3$ を代入したもので、
この形の行列は多項式フィッティング(第I部)にも現れた。
\end{rei}

\begin{rei}[例題・標準(和空間への分解)]\label{rei:sub-decomp}
$W_1 = \langle (1, 0, 1) \rangle$,
$W_2 = \{\, (x, y, z) \mid x + y + z = 0 \,\}$ について、
$W_1 + W_2 = \R^3$ かつ $W_1 \cap W_2 = \{\vect{0}\}$ を示せ。\\
\textbf{解.}\ 任意の $\vect{v} = (x, y, z)$ を
$\vect{v} = t\,(1, 0, 1) + \vect{w}$($\vect{w} \in W_2$)と
分解できるかを考える。$\vect{w} = \vect{v} - t\,(1, 0, 1)
= (x - t,\ y,\ z - t)$ が $W_2$ に入る条件は
\[
(x - t) + y + (z - t) = 0, \qquad \text{すなわち} \qquad
t = \frac{x + y + z}{2}.
\]
この $t$ を選べば必ず分解できるから $W_1 + W_2 = \R^3$ である。
たとえば $\vect{v} = (3, 1, 2)$ なら $t = 3$ で
$(3,1,2) = 3\,(1,0,1) + (0, 1, -1)$($0 + 1 - 1 = 0$ を満たす)。

共通部分:$t\,(1, 0, 1) \in W_2$ とすると $t + 0 + t = 2t = 0$ より
$t = 0$。よって $W_1 \cap W_2 = \{\vect{0}\}$ である。
この場合、分解は一意になる(二通りに書けたとすると、差をとれば
$W_1 \cap W_2$ の零でない元ができて矛盾)。共通部分が $\{\vect{0}\}$ の
和空間はとくに\textbf{直和}と呼ばれ、第V部以降で活躍する。
\end{rei}

\begin{matome}
\begin{itemize}
\item \textbf{部分空間}とは、和とスカラー倍で閉じた空でない部分集合である。
閉じてさえいれば、$8$ つの公理は親の空間から遺伝して、
それ自身ベクトル空間になる(定理 \ref{thm:sub-inherit})。
\item 判定の第一歩は「$\vect{0}$ を含むか」。含まなければ即座に否定できる。
ただし含んでいても閉性が破れる集合はある(例題 \ref{rei:sub-curve})。
破れ方は和・スカラー倍のどちらでも起こりうる
(反例 \ref{hanrei:sub-union}・\ref{hanrei:sub-lattice})。
\item \textbf{共通部分} $W_1 \cap W_2$ と\textbf{和空間} $W_1 + W_2$ は
部分空間になる。\textbf{和集合} $W_1 \cup W_2$ はそうとは限らず、
それを含む最小の部分空間が $W_1 + W_2$ である。
\item $\langle \vect{a}_1, \dots, \vect{a}_k \rangle$(線形結合全体)は、
$\vect{a}_1, \dots, \vect{a}_k$ を含む\textbf{最小の部分空間}である
(定理 \ref{thm:span-min})。
\item $\R^n$ では、「$\vect{b}$ は張られた空間に入るか」は連立方程式の
可解性($\mathrm{rank}\,A = \mathrm{rank}\,(A \mid \vect{b})$)に、
「$\R^n$ 全体を張るか」は $\mathrm{rank}\,A = n$(正方なら
$\det A \neq 0$)に帰着する。
\end{itemize}
\end{matome}

\section*{章末問題}

\begin{mondai}[基本]\label{q:sub-ex1}
次の集合はそれぞれ部分空間か。部分空間であるものは証明し、
そうでないものは理由を一つ挙げよ。\\
　(1) $\{\, (x, y, z) \in \R^3 \mid x - 3y + 2z = 0 \,\}$\\
　(2) $\{\, (x, y) \in \R^2 \mid y = x^2 \,\}$\\
　(3) $\{\, (x, y) \in \R^2 \mid x \le 0 \,\}$
\end{mondai}

\begin{mondai}[基本]\label{q:sub-ex2}
$\vect{a}_1 = \mat{1 \\ 0 \\ 2}$, $\vect{a}_2 = \mat{1 \\ 1 \\ 0}$ について、
$\vect{b} = \mat{3 \\ 1 \\ 4}$ と $\vect{b}' = \mat{1 \\ 1 \\ 1}$ が
$\langle \vect{a}_1, \vect{a}_2 \rangle$ に属するか判定し、
属する場合は線形結合の形に書け。
\end{mondai}

\begin{mondai}[基本]\label{q:sub-ex3}
$2 \times 2$ の対角行列全体
$D = \{\, \mat{a & 0 \\ 0 & d} \mid a, d \in \R \,\}$
が $M_{2,2}$ の部分空間であることを示せ。
\end{mondai}

\begin{mondai}[基本]\label{q:sub-ex4}
$\R[x]_2$ の部分集合 $W = \{\, p \in \R[x]_2 \mid p(0) = 0 \,\}$ が
部分空間であることを示せ。$W$ の元はどのような形の多項式か。
\end{mondai}

\begin{mondai}[標準]\label{q:sub-ex5}
$\R^2$ の二つの部分空間
$W_1 = \{\, (x, y) \mid x + y = 0 \,\}$,
$W_2 = \{\, (x, y) \mid x - y = 0 \,\}$ について、次を示せ。\\
　(1) 和集合 $W_1 \cup W_2$ は部分空間でない。\\
　(2) 和空間について $W_1 + W_2 = \R^2$ が成り立つ。
\end{mondai}

\begin{mondai}[標準]\label{q:sub-ex6}
$\R[x]_2$ の部分集合
$U_1 = \{\, p \mid p(1) = 0 \,\}$ と $U_2 = \{\, p \mid p(1) = 1 \,\}$
について、$U_1$ は部分空間であり $U_2$ は部分空間でないことを示せ。
\end{mondai}

\begin{mondai}[標準]\label{q:sub-ex7}
$\R^3$ の二つの平面
$W_1 = \{\, (x, y, z) \mid x + y - z = 0 \,\}$,
$W_2 = \{\, (x, y, z) \mid 2x - y + z = 0 \,\}$ の共通部分
$W_1 \cap W_2$ を、$\langle \vect{a} \rangle$(あるベクトルの張る直線)の
形で表せ。
\end{mondai}

\begin{mondai}[標準]\label{q:sub-ex8}
次の生成された部分空間はそれぞれ何か($\R^2$ の直線か、$\R^2$ 全体か)。
理由とともに答えよ。\\
　(1) $\left\langle \mat{1 \\ 2},\ \mat{2 \\ 4} \right\rangle$\\
　(2) $\left\langle \mat{1 \\ 2},\ \mat{1 \\ 3} \right\rangle$
\end{mondai}

\begin{mondai}[発展]\label{q:sub-ex9}
$\vect{a}_1 = \mat{1 \\ 0 \\ 1}$, $\vect{a}_2 = \mat{0 \\ 1 \\ 1}$,
$\vect{a}_3 = \mat{1 \\ 1 \\ 2}$ は $\R^3$ を生成するか。
生成しない場合は、$\langle \vect{a}_1, \vect{a}_2, \vect{a}_3 \rangle$ が
どのような平面か、$x, y, z$ の方程式で表せ。
\end{mondai}

\begin{mondai}[発展]\label{q:sub-ex10}
$C(\R)$ の部分空間である偶関数全体 $E$(問 \ref{q:vs-ex7})と
奇関数全体 $O = \{\, f \mid f(-x) = -f(x) \,\}$ について、次を示せ。\\
　(1) $E \cap O = \{o\}$($o$ は零関数)。\\
　(2) $E + O = C(\R)$。すなわち、どんな連続関数も
偶関数と奇関数の和に分解できる。
\end{mondai}

\begin{mondai}[証明]\label{q:sub-ex11}
$\vect{b} \in \langle \vect{a}_1, \dots, \vect{a}_k \rangle$ ならば
\[
\langle \vect{a}_1, \dots, \vect{a}_k, \vect{b} \rangle
= \langle \vect{a}_1, \dots, \vect{a}_k \rangle
\]
であることを証明せよ(すでに張られている空間に入っているベクトルを
生成系に付け加えても、張られる空間は変わらない)。
\end{mondai}

\begin{mondai}[証明]\label{q:sub-ex12}
$W_1, W_2$ を $V$ の部分空間とする。和集合 $W_1 \cup W_2$ が
部分空間になるのは、$W_1 \subseteq W_2$ または $W_2 \subseteq W_1$ の
ときに限ることを証明せよ。
\end{mondai}

\bigskip
{\small
\noindent\textbf{【章末問題の方針】}\quad
まず自力で取り組み、行き詰まったら下の方針を見よ。記述による完全解答は
巻末「問の解答」にある。
\begin{itemize}
\item \textbf{問 \ref{q:sub-ex1}}:まず原点を含むか。含むものは
文字で閉性を確かめる。(2) はスカラー倍、(3) は負のスカラー倍を疑う。
\item \textbf{問 \ref{q:sub-ex2}}:$c_1\vect{a}_1 + c_2\vect{a}_2 = \vect{b}$
を成分ごとの連立方程式にして解く。矛盾が出れば属さない。
\item \textbf{問 \ref{q:sub-ex3}}:対角行列どうしの和・スカラー倍が
対角行列のままであることを成分で確かめる。
\item \textbf{問 \ref{q:sub-ex4}}:$p(0)$ は定数項である。
$(p + q)(0)$ と $(cp)(0)$ を計算する。
\item \textbf{問 \ref{q:sub-ex5}}:(1) は各直線から $1$ 点ずつとって
和を作る。(2) は任意の $(x, y)$ を二つの直線上のベクトルの和に
実際に分解する(係数を未知数にして解く)。
\item \textbf{問 \ref{q:sub-ex6}}:$U_2$ は零多項式を含むか。
$U_1$ は $(p+q)(1)$, $(cp)(1)$ を計算する。
\item \textbf{問 \ref{q:sub-ex7}}:二つの方程式を連立して掃き出す。
自由変数が $1$ 個残り、直線になる。
\item \textbf{問 \ref{q:sub-ex8}}:$2$ 本が平行(一方が他方の定数倍)か
どうかを見る。平行でなければ行列式で判定する。
\item \textbf{問 \ref{q:sub-ex9}}:行列式を計算する。$0$ なら
従属な $1$ 本を見つけて捨て、残り $2$ 本の張る平面の法線を求める。
\item \textbf{問 \ref{q:sub-ex10}}:(1) は $f(x) = f(-x) = -f(x)$ から。
(2) は $\dfrac{f(x) + f(-x)}{2}$ と $\dfrac{f(x) - f(-x)}{2}$ を考える。
\item \textbf{問 \ref{q:sub-ex11}}:相互の包含を示す。$\supseteq$ は
係数 $0$ を付け足すだけ。$\subseteq$ は $\vect{b}$ を線形結合で
書き直して整理する。
\item \textbf{問 \ref{q:sub-ex12}}:一方が他方を含めば和集合は大きい方
そのもの。どちらも含まないときは、$\vect{u} \in W_1 \setminus W_2$ と
$\vect{v} \in W_2 \setminus W_1$ の和がどちらにも入れないことを示す。
\end{itemize}
}

% =====================================================
\chapter{基底と次元}\label{chap:basis}
% =====================================================

\begin{goal}
\begin{itemize}
\item 基底の二つの条件(線形独立+生成)と\textbf{座標}の意味を理解し、
与えられた組が基底かどうかの判定と、基底に関する座標の計算が
できるようになる。
\item 基底の本数が選び方によらないこと(\textbf{次元}が定義できること)の
証明を、交換補題から順を追って理解する。
\item $n$ 次元空間では「独立な $n$ 本」または「生成する $n$ 本」の
\textbf{どちらか一方の確認だけで}基底と判定できることを使いこなせるようになる。
\end{itemize}
\end{goal}

\section{動機 ― 座標軸は自分で選ぶもの}

$\R^n$ のベクトルには生まれつき成分が並んでいる。しかし前章までに
仲間入りした空間たち――多項式の空間、解空間、漸化式を満たす数列の空間
――には、あらかじめ決まった「座標軸」はない。
それでも、たとえば平面 $x + y + z = 0$ 上の点は $2$ 個のパラメータで
過不足なく表せた(第II部)。軸は与えられていなくても、
\textbf{自分で軸を選べば座標が引ける}のである。

この「自分で選ぶ座標軸」を厳密にしたものが\textbf{基底}であり、
軸の本数が\textbf{次元}である。ただし一つ、確かめるべきことがある。
座標軸の選び方は無数にあるのだから、
\textbf{選び方によって軸の本数が変わってしまわないか?}
もし $3$ 本の基底と $4$ 本の基底が同じ空間に共存できたら、
「次元」という言葉は意味を失う。本章の山場は、
その心配が無用であることの証明(基底の本数の不変性)である。

\section{基底 ― 空間の「座標軸」}

線形独立の定義(第II部の定義 \ref{def:linindep})は、
線形結合だけで書かれていたから、一般のベクトル空間でもそのまま通用する。
念のため書き直しておく。

\begin{teigi}[線形独立・線形従属(一般のベクトル空間)]\label{def:vs-indep}
ベクトル空間 $V$ のベクトル $\vect{v}_1, \dots, \vect{v}_k$ が
\[
c_1\vect{v}_1 + \cdots + c_k\vect{v}_k = \vect{0}
\implies c_1 = \cdots = c_k = 0
\]
を満たすとき\textbf{線形独立}であるといい、
そうでないとき(係数のどれかが $0$ でない線形結合で $\vect{0}$ が
作れるとき)\textbf{線形従属}であるという。
\end{teigi}

\begin{teigi}[基底]
ベクトル空間 $V$ のベクトルの組 $\vect{e}_1, \dots, \vect{e}_n$ が
\begin{enumerate}
\item 線形独立であり、
\item $V$ を生成する($V$ のすべてのベクトルがその線形結合で書ける)
\end{enumerate}
とき、この組を $V$ の\textbf{基底}という。
\end{teigi}

\begin{teiri}[座標の一意性]\label{thm:coord-unique}
$\vect{e}_1, \dots, \vect{e}_n$ が $V$ の基底ならば、
任意の $\vect{v} \in V$ は
\[
\vect{v} = c_1 \vect{e}_1 + \cdots + c_n \vect{e}_n
\]
の形に\textbf{ただ一通り}に書ける。この $(c_1, \dots, c_n)$ を、
この基底に関する $\vect{v}$ の\textbf{座標}という。
\end{teiri}

\noindent\textbf{(証明)}\quad
書ける(存在):基底の条件2(生成)そのものである。
一通り(一意):二通りに
\[
\vect{v} = c_1 \vect{e}_1 + \cdots + c_n \vect{e}_n
= d_1 \vect{e}_1 + \cdots + d_n \vect{e}_n
\]
と書けたとして、辺々引くと
\[
(c_1 - d_1)\vect{e}_1 + \cdots + (c_n - d_n)\vect{e}_n = \vect{0}.
\]
基底の条件1(独立)より係数はすべて $0$、すなわち $c_i = d_i$ である。
$\blacksquare$

つまり基底とは、\textbf{空間のすべての点に過不足なく「住所」を与える
座標軸の組}である。生成が「すべての点に住所が行き渡ること」を、
独立が「住所の重複がないこと」を保証している。

\begin{rei}
\begin{itemize}
\item $\R^n$ の\textbf{標準基底}:
$\vect{e}_1 = \mat{1 \\ 0 \\ \vdots \\ 0}, \dots,
\vect{e}_n = \mat{0 \\ \vdots \\ 0 \\ 1}$。
\item $\R[x]_n$ の基底:$1, x, x^2, \dots, x^n$。
多項式 $a_0 + a_1 x + \cdots + a_n x^n$ の係数の並びがそのまま座標である。
\item 基底は一通りではない。$\R^2$ では
$\mat{1 \\ 1}$ と $\mat{1 \\ -1}$ も基底になる(独立で、生成することを確かめよ)。
\textbf{どの基底を選ぶかは人間の自由}であり、この自由こそ第IV部以降の主題
(うまい基底選び=対角化)につながる。
\end{itemize}
\end{rei}

\section{次元 ― 本数の不変性を証明する}

動機で述べた核心の問いに答えよう。鍵は次の補題である。

\begin{teiri}[交換補題]\label{thm:exchange}
ベクトル空間 $V$ が $n$ 本のベクトル $\vect{u}_1, \dots, \vect{u}_n$ で
生成されるならば、$V$ の中のどんな $n + 1$ 本のベクトルも線形従属である
($n + 2$ 本以上でも同様)。
\end{teiri}

\noindent\textbf{(証明)}\quad
$\vect{v}_1, \dots, \vect{v}_{n+1} \in V$ を任意にとる。
生成の仮定により、各 $\vect{v}_j$ は
\[
\vect{v}_j = a_{1j}\vect{u}_1 + a_{2j}\vect{u}_2 + \cdots + a_{nj}\vect{u}_n
\qquad (j = 1, \dots, n+1)
\]
と書ける。すべては $0$ でない係数 $c_1, \dots, c_{n+1}$ で
$c_1\vect{v}_1 + \cdots + c_{n+1}\vect{v}_{n+1} = \vect{0}$ とできることを
示せばよい。上の表示を代入し、$\vect{u}_i$ ごとにまとめると
\[
c_1\vect{v}_1 + \cdots + c_{n+1}\vect{v}_{n+1}
= \sum_{i=1}^{n} \bigl( a_{i1}c_1 + a_{i2}c_2 + \cdots + a_{i,n+1}c_{n+1} \bigr)
\vect{u}_i
\]
となる。そこで、$c_1, \dots, c_{n+1}$ に関する同次連立一次方程式
\[
a_{i1}c_1 + a_{i2}c_2 + \cdots + a_{i,n+1}c_{n+1} = 0
\qquad (i = 1, \dots, n)
\]
を考えると、これは\textbf{式 $n$ 本・未知数 $n + 1$ 個}の同次方程式だから、
非自明な解 $(c_1, \dots, c_{n+1}) \neq (0, \dots, 0)$ をもつ
(第II部の定理 \ref{thm:wide-homog})。この $c_j$ を使えば、
各 $\vect{u}_i$ の係数がすべて $0$ になるので
$c_1\vect{v}_1 + \cdots + c_{n+1}\vect{v}_{n+1} = \vect{0}$。
よって $\vect{v}_1, \dots, \vect{v}_{n+1}$ は線形従属である。
($n + 2$ 本以上の場合は、余分なベクトルの係数を $0$ とおけば
$n + 1$ 本の場合に帰着する。) $\blacksquare$

抽象的な「次元」の正当性が、掃き出し法の具体的な定理
(未知数が式より多い同次方程式には非自明解がある)に支えられている
――抽象と計算の往復が線形代数の醍醐味である。

\begin{teiri}[基底の本数の不変性]\label{thm:dim-invariant}
ベクトル空間 $V$ が $n$ 本のベクトルからなる基底をもつならば、
$V$ のどの基底もちょうど $n$ 本からなる。
\end{teiri}

\noindent\textbf{(証明)}\quad
$m$ 本の基底 $B$ と $n$ 本の基底 $B'$ があったとして、$m < n$ と仮定する。
$B$ は $V$ を $m$ 本で生成するから、交換補題 \ref{thm:exchange} により
$V$ の中の $m + 1$ 本以上のベクトルは必ず従属である。
ところが $B'$ は $n \ge m + 1$ 本の\textbf{線形独立な}組であり、矛盾。
$n < m$ と仮定しても同様に矛盾するから、$m = n$ である。 $\blacksquare$

\begin{teigi}[次元]
この $n$ を $V$ の\textbf{次元}といい、$\dim V$ と書く。
有限本の基底をもつ空間を\textbf{有限次元}、もたない空間を\textbf{無限次元}という。
\end{teigi}

\begin{rei}
$\dim \R^n = n$、$\dim \R[x]_n = n + 1$($1, x, \dots, x^n$ の $n+1$ 本)。
同次方程式 $A\vect{x} = \vect{0}$ の解空間の次元は、
第II部の自由度 $n - \mathrm{rank}\,A$ に等しい
(パラメータ表示に現れる $n - \mathrm{rank}\,A$ 本のベクトルが基底になる)。
一方、$\R[x]$ や $C(\R)$ は無限次元である
($1, x, x^2, \dots$ は何本とっても線形独立)。
\end{rei}

\section{基底の作り方と $n$ 本判定}

基底は「間引く」か「付け足す」かで必ず作れる。

\begin{teiri}[間引きと拡張]\label{thm:dim-extend}
$V$ を有限次元ベクトル空間とする。
\begin{enumerate}
\item \textbf{(間引き)}\ $V$ の有限な生成系からは、一部のベクトルを
取り除いて $V$ の基底を選び出せる。
\item \textbf{(拡張)}\ $V$ の線形独立な組には、ベクトルを付け足して
$V$ の基底に拡張できる。
\end{enumerate}
\end{teiri}

\noindent\textbf{(証明)}\quad
(1) 生成系が従属なら、あるベクトルが残りの線形結合で書ける
(第II部の定理 \ref{thm:dep-combo} と同じ論法)。そのベクトルを
取り除いても張る空間は変わらない(問 \ref{q:sub-ex11})。
これを繰り返すと、生成系は有限本だからいつか止まり、
残った組は独立な生成系、すなわち基底である。

(2) 鍵は次の補題である:
\textbf{独立な組 $\vect{u}_1, \dots, \vect{u}_k$ と、その張る空間に
入らないベクトル $\vect{v}$ をあわせた $k + 1$ 本も独立である。}
実際 $c_1\vect{u}_1 + \cdots + c_k\vect{u}_k + c\vect{v} = \vect{0}$ とすると、
もし $c \neq 0$ なら $\vect{v} = -\dfrac{c_1}{c}\vect{u}_1 - \cdots
- \dfrac{c_k}{c}\vect{u}_k$ となって
$\vect{v} \in \langle \vect{u}_1, \dots, \vect{u}_k \rangle$ に矛盾。
よって $c = 0$ で、残りは $\vect{u}_i$ の独立性より $c_i = 0$ となる。

そこで、与えられた独立な組がまだ $V$ を生成していなければ、
張る空間の外からベクトルを $1$ 本とって付け足す。補題により独立性は
保たれる。$\dim V = n$ とすると、交換補題 \ref{thm:exchange} により
独立な組は $n$ 本以下だから、この操作は必ず止まる。
止まったとき、その組は「張る空間の外にベクトルがない」、
すなわち $V$ を生成する独立系=基底である。 $\blacksquare$

\begin{teiri}[$n$ 次元空間での基底判定]\label{thm:dim-ncheck}
$\dim V = n$ とする。$V$ のちょうど $n$ 本のベクトル
$\vect{v}_1, \dots, \vect{v}_n$ について、
\begin{enumerate}
\item 線形独立ならば、基底である(生成の確認は不要)。
\item $V$ を生成するならば、基底である(独立の確認は不要)。
\end{enumerate}
\end{teiri}

\noindent\textbf{(証明)}\quad
(1) 拡張定理 \ref{thm:dim-extend}(2) により、この $n$ 本は基底に
拡張できる。しかし基底は必ず $n$ 本(定理 \ref{thm:dim-invariant})
だから、付け足すべきベクトルは $0$ 本であり、この組がすでに基底である。
(2) 間引き定理 \ref{thm:dim-extend}(1) により、この $n$ 本から基底を
選び出せる。基底は $n$ 本だから、取り除くべきベクトルは $0$ 本であり、
この組がすでに基底である。 $\blacksquare$

次元が分かっている空間では、この定理のおかげで基底の判定が半分の手間で
済む。たとえば $\R^n$ の $n$ 本なら、行列式が $0$ でないこと
(=独立)を見るだけでよい。

\begin{teiri}[部分空間の次元]\label{thm:dim-sub}
$n$ 次元ベクトル空間 $V$ の部分空間 $W$ について、次が成り立つ。
\begin{enumerate}
\item $W$ も有限次元で、$\dim W \le n$。
\item $\dim W = n$ ならば $W = V$。
\end{enumerate}
\end{teiri}

\noindent\textbf{(証明)}\quad
(1) $W$ の中の独立な組は $V$ の中でも独立だから、
交換補題 \ref{thm:exchange} により本数は $n$ 以下である。
$W$ の中で独立な組を、拡張定理の証明と同じ要領で
($W$ の中の、張る空間の外のベクトルを付け足して)延ばしていくと、
$n$ 本以下で必ず止まる。止まった組は $W$ を生成する独立系、
すなわち $W$ の基底だから、$\dim W \le n$。
(2) $\dim W = n$ なら、$W$ の基底は $V$ の中の独立な $n$ 本でもあるから、
定理 \ref{thm:dim-ncheck}(1) により $V$ の基底である。よって
$V$ のすべてのベクトルがその線形結合、すなわち $W$ の元であり、
$W = V$ となる。 $\blacksquare$

\begin{mondai}\label{q:basis1}
　(1) 平面 $W = \{\,(x, y, z) \in \R^3 \mid x - y + 2z = 0\,\}$ の基底を一組求め、
$\dim W$ を答えよ。\\
　(2) $1,\ 1 + x,\ (1 + x)^2$ が $\R[x]_2$ の基底であることを示せ。
\end{mondai}

\section{反例 ― ありがちな誤解を正す}

\begin{hanrei}[本数が合っていても基底とは限らない]\label{hanrei:dim-count}
$\R^2$ で $\mat{1 \\ 2}, \mat{2 \\ 4}$ は、本数こそ $\dim \R^2 = 2$ に
一致するが、$\mat{2 \\ 4} = 2\mat{1 \\ 2}$ と従属なので基底ではない
(張るのは直線 $y = 2x$ だけで、生成もしていない)。
逆に、独立でも本数が足りなければ基底ではない:
$\mat{1 \\ 0 \\ 0}, \mat{0 \\ 1 \\ 0}$ は $\R^3$ で独立だが、
$\mat{0 \\ 0 \\ 1}$ を書き表せないので生成しない。
定理 \ref{thm:dim-ncheck} が使えるのは、
\textbf{本数がちょうど $\dim V$ に一致しているとき}だけであり、
そのときも独立か生成のどちらか一方は確認しなければならない。
\end{hanrei}

\begin{hanrei}[「座標」は基底とセットでしか意味をもたない]\label{hanrei:dim-coord}
$\R^2$ のベクトル $\vect{v} = \mat{3 \\ 1}$ の座標は $(3, 1)$
――と言い切るのは早計である。それは\textbf{標準基底に関する}座標にすぎない。
基底 $\vect{e}'_1 = \mat{1 \\ 1}$, $\vect{e}'_2 = \mat{1 \\ -1}$ をとると
\[
\mat{3 \\ 1} = 2\mat{1 \\ 1} + 1\mat{1 \\ -1}
\]
だから、この基底に関する $\vect{v}$ の座標は $(2, 1)$ である。
同じベクトルでも、基底(=座標軸)が変われば住所は変わる。
「座標」を口にするときは、どの基底に関する座標かをつねに意識せよ
(この観点は第 \ref{chap:repmat} 章の主題になる)。
\end{hanrei}

\section{例題}

基底の判定・座標の計算・次元の決定を練習する。
以下は\textbf{解答つきの例題}である。手を動かして追ってほしい。

\begin{rei}[例題・基本(解空間の基底)]\label{rei:dim-plane}
平面 $W = \{\,(x, y, z) \in \R^3 \mid x + 2y - z = 0\,\}$ の基底を
一組求め、$\dim W$ を答えよ。\\
\textbf{解.}\ $x = -2y + z$ より、$y = s$, $z = t$ とおくと
\[
\mat{x \\ y \\ z} = s\mat{-2 \\ 1 \\ 0} + t\mat{1 \\ 0 \\ 1}.
\]
この $2$ 本は $W$ を生成し、しかも独立である(第 $2$・第 $3$ 成分を
見れば、$s\,(-2,1,0) + t\,(1,0,1) = \vect{0}$ から $s = 0$, $t = 0$ が
直ちに出る)。よって基底をなし、$\dim W = 2$ である。
パラメータ表示では「自由変数の位置」に $1$ と $0$ が振り分けられるため、
この方法で得られる生成系は自動的に独立になる。
\end{rei}

\begin{rei}[例題・基本(基底の判定と座標)]\label{rei:dim-coord}
$\vect{a}_1 = \mat{1 \\ 1}$, $\vect{a}_2 = \mat{1 \\ -1}$ が $\R^2$ の
基底であることを示し、$\vect{v} = \mat{1 \\ 5}$ のこの基底に関する
座標を求めよ。\\
\textbf{解.}\ $\det \mat{1 & 1 \\ 1 & -1} = -1 - 1 = -2 \neq 0$ より
$2$ 本は独立(第II部の定理 \ref{thm:reg-rank})。
$\dim \R^2 = 2$ だから、定理 \ref{thm:dim-ncheck}(1) により基底である。
座標は $c_1\vect{a}_1 + c_2\vect{a}_2 = \vect{v}$ を解いて求める:
\[
c_1 + c_2 = 1, \qquad c_1 - c_2 = 5
\]
より $c_1 = 3$, $c_2 = -2$。座標は $(3, -2)$ である
(検算:$3\mat{1\\1} - 2\mat{1\\-1} = \mat{1\\5}$)。
\end{rei}

\begin{rei}[例題・基本(行列の空間の次元)]\label{rei:dim-sym}
$2 \times 2$ 対称行列全体 $S = \{\, A \mid {}^t\!A = A \,\}$
(問 \ref{q:vs-ex9})の基底を一組求め、$\dim S$ を答えよ。\\
\textbf{解.}\ 対称行列は $\mat{a & b \\ b & c}$ の形だから、
\[
\mat{a & b \\ b & c}
= a\mat{1 & 0 \\ 0 & 0} + b\mat{0 & 1 \\ 1 & 0} + c\mat{0 & 0 \\ 0 & 1}
\]
と、$3$ 枚の行列の線形結合で書ける($3$ 枚は $S$ を生成する)。
また線形結合が $O$ になるのは、各成分を比べて $a = b = c = 0$ の
ときだけだから独立でもある。よって基底をなし、$\dim S = 3$ である。
なお $\dim M_{2,2} = 4$(成分一つだけが $1$ の「行列単位」$4$ 枚が基底)
だから、$S$ は $4$ 次元空間の中の $3$ 次元部分空間である。
\end{rei}

\begin{rei}[例題・標準(多項式の基底と座標)]\label{rei:dim-poly}
$1,\ 1 + x,\ 1 + x + x^2$ が $\R[x]_2$ の基底であることを示し、
$p(x) = 2 + 3x + x^2$ のこの基底に関する座標を求めよ。\\
\textbf{解.}\ $c_1 \cdot 1 + c_2(1 + x) + c_3(1 + x + x^2) = 0$
(恒等的に)とすると、次数ごとに係数を比べて
\[
c_1 + c_2 + c_3 = 0, \qquad c_2 + c_3 = 0, \qquad c_3 = 0
\]
となり、下から順に $c_3 = 0$, $c_2 = 0$, $c_1 = 0$。よって独立である。
$\dim \R[x]_2 = 3$ だから、定理 \ref{thm:dim-ncheck}(1) により基底である。
座標は $c_1 + c_2(1 + x) + c_3(1 + x + x^2) = 2 + 3x + x^2$ の係数比較で
\[
c_1 + c_2 + c_3 = 2, \qquad c_2 + c_3 = 3, \qquad c_3 = 1
\]
より $c_3 = 1$, $c_2 = 2$, $c_1 = -1$。座標は $(-1, 2, 1)$ である
(検算:$-1 + 2(1 + x) + (1 + x + x^2) = 2 + 3x + x^2$)。
\end{rei}

\begin{rei}[例題・標準(生成系から基底を間引く)]\label{rei:dim-select}
$\vect{a}_1 = \mat{1 \\ 2 \\ 3}$, $\vect{a}_2 = \mat{2 \\ 4 \\ 6}$,
$\vect{a}_3 = \mat{1 \\ 1 \\ 1}$, $\vect{a}_4 = \mat{2 \\ 3 \\ 4}$ の
張る部分空間 $W = \langle \vect{a}_1, \vect{a}_2, \vect{a}_3, \vect{a}_4
\rangle$ の基底をこの中から選び出し、$\dim W$ を答えよ。\\
\textbf{解.}\ 間引き定理 \ref{thm:dim-extend}(1) の方針で、従属なものを
順に捨てる。$\vect{a}_2 = 2\vect{a}_1$ だから $\vect{a}_2$ を捨ててよい
(問 \ref{q:sub-ex11}:張る空間は変わらない)。
$\vect{a}_1, \vect{a}_3$ は比例していないので独立。
$\vect{a}_4$ は $\vect{a}_4 = \vect{a}_1 + \vect{a}_3$
(成分ごとに $2 = 1 + 1$, $3 = 2 + 1$, $4 = 3 + 1$)と書けるから捨ててよい。
残った $\vect{a}_1, \vect{a}_3$ は独立な生成系なので $W$ の基底であり、
$\dim W = 2$($W$ は $\R^3$ 内の平面)である。
\end{rei}

\begin{rei}[例題・標準(数列の空間の基底)]\label{rei:dim-seq}
漸化式 $a_{n+2} = 2a_{n+1} - a_n$ を満たす数列全体 $W$
(問 \ref{q:vs-ex6})について、定数数列 $(1, 1, 1, \dots)$ と
数列 $(0, 1, 2, 3, \dots)$($a_n = n$)が $W$ の基底であることを示せ。\\
\textbf{解.}\ まず両者が $W$ に属すること:定数数列は
$1 = 2 \cdot 1 - 1$、数列 $a_n = n$ は $n + 2 = 2(n + 1) - n$ より、
どちらも漸化式を満たす。

生成:任意の $(a_n) \in W$ に対し、二つの数列の線形結合
$b_n = a_0 \cdot 1 + (a_1 - a_0) \cdot n$ で定まる数列 $(b_n)$ は
$W$ に属し、$b_0 = a_0$, $b_1 = a_1$ と最初の $2$ 項が $(a_n)$ と一致する。
漸化式は最初の $2$ 項から残りすべてを決めるので $(a_n) = (b_n)$、
すなわち $(a_n)$ は $2$ つの数列の線形結合である。

独立:$c_1 \cdot 1 + c_2 \cdot n = 0$(すべての $n$)とすると、
$n = 0$ で $c_1 = 0$、$n = 1$ で $c_2 = 0$。

よって基底をなし、$\dim W = 2$ である。等差数列の一般項
$a_n = a_0 + (a_1 - a_0)n$ は、この基底に関する座標表示だったのである。
\end{rei}

\begin{itembox}[l]{深掘り:無限次元の世界}
$\R[x]$(次数制限なしの多項式全体)には有限本の基底がない。
どんな有限個の多項式をとっても、その中の最高次数を $d$ とすれば
$x^{d+1}$ はそれらの線形結合で書けないからである。
$1, x, x^2, \dots$ は何本とっても独立で、$\R[x]$ は無限次元である。
連続関数の空間 $C(\R)$ はさらに広大で、$e^x$ や $\sin x$ は
どんな多項式とも一致しない(多項式は有限回微分すれば $0$ になるが、
$e^x$ は何回微分しても $e^x$ のままである)。

無限次元空間で「基底」を扱うには繊細さが要る。線形結合とは
あくまで\textbf{有限個}の和だから、
$e^x = 1 + x + \frac{x^2}{2!} + \cdots$ のような\textbf{無限和}は、
ベクトル空間の言葉のままでは扱えない。無限和に意味を与えるには
「収束」の概念、すなわち解析学との合流が必要で、その先に広がるのが
\textbf{関数解析}(フーリエ級数、量子力学の数学)である。
有限次元の理論はその登山口であり、本章で証明した「次元の不変性」の
ような足場固めは、無限次元でも議論の模範になる。
\end{itembox}

\begin{mondai}\label{q:basis2}
問 \ref{q:fib} の空間 $V$(漸化式 $a_{n+2} = a_{n+1} + a_n$ を満たす数列全体)
について考える。数列 $(a_n) \in V$ は最初の $2$ 項 $a_0, a_1$ を決めれば
すべて決まることに注意し、$\dim V = 2$ であることを示せ。
(ヒント:$a_0 = 1, a_1 = 0$ で始まる数列と $a_0 = 0, a_1 = 1$ で始まる数列が
基底になる。)
\end{mondai}

\begin{matome}
\begin{itemize}
\item \textbf{基底}=線形独立な生成系。生成が「すべての点に住所を配る」こと、
独立が「住所の重複を防ぐ」ことを保証し、座標が\textbf{ただ一通り}に定まる
(定理 \ref{thm:coord-unique})。座標は基底とセットでのみ意味をもつ
(反例 \ref{hanrei:dim-coord})。
\item \textbf{交換補題}:$n$ 本で生成される空間では $n + 1$ 本以上は必ず従属。
証明は「未知数が式より多い同次方程式には非自明解がある」に帰着する。
これにより\textbf{基底の本数はとり方によらず}(定理 \ref{thm:dim-invariant})、
\textbf{次元} $\dim V$ が矛盾なく定義できる。
\item 基底は、生成系を\textbf{間引く}か、独立系に\textbf{付け足す}かで
必ず作れる(定理 \ref{thm:dim-extend})。
\item $\dim V = n$ のとき、$n$ 本のベクトルは「独立」か「生成」の
\textbf{どちらか一方}を確かめれば基底である(定理 \ref{thm:dim-ncheck})。
本数が $n$ でなければこの近道は使えない(反例 \ref{hanrei:dim-count})。
\item 部分空間 $W$ は $\dim W \le \dim V$ を満たし、
$\dim W = \dim V$ なら $W = V$(定理 \ref{thm:dim-sub})。
\item 解空間の基底はパラメータ表示から読み取れ、次元は自由度
$n - \mathrm{rank}\,A$ に一致する。
\end{itemize}
\end{matome}

\section*{章末問題}

\begin{mondai}[基本]\label{q:dim-ex1}
平面 $W = \{\,(x, y, z) \in \R^3 \mid 2x - y + 3z = 0\,\}$
(例題 \ref{rei:sub-plane} の部分空間)の基底を一組求め、$\dim W$ を答えよ。
\end{mondai}

\begin{mondai}[基本]\label{q:dim-ex2}
$\vect{a}_1 = \mat{2 \\ 1}$, $\vect{a}_2 = \mat{1 \\ 1}$ が $\R^2$ の
基底であることを示し、$\vect{v} = \mat{1 \\ 4}$ のこの基底に関する座標を求めよ。
\end{mondai}

\begin{mondai}[基本]\label{q:dim-ex3}
$2 \times 2$ の対角行列全体 $D$(問 \ref{q:sub-ex3})の基底を一組求め、
$\dim D$ を答えよ。
\end{mondai}

\begin{mondai}[基本]\label{q:dim-ex4}
$W = \{\, p \in \R[x]_2 \mid p(0) = 0 \,\}$(問 \ref{q:sub-ex4})の
基底を一組求め、$\dim W$ を答えよ。
\end{mondai}

\begin{mondai}[標準]\label{q:dim-ex5}
$\R^4$ の部分空間
$W = \{\,(x_1, x_2, x_3, x_4) \mid x_1 + x_2 + x_3 + x_4 = 0\,\}$ の
基底を一組求め、$\dim W$ を答えよ。
\end{mondai}

\begin{mondai}[標準]\label{q:dim-ex6}
$1 + x,\ 1 - x,\ x^2$ が $\R[x]_2$ の基底であることを示し、
$p(x) = 3 + x - 2x^2$ のこの基底に関する座標を求めよ。
\end{mondai}

\begin{mondai}[標準]\label{q:dim-ex7}
$2 \times 2$ の交代行列全体 $T = \{\, A \mid {}^t\!A = -A \,\}$
(問 \ref{q:vs-ex9})の基底を一組求め、$\dim T$ を答えよ。
また、$\dim S + \dim T$($S$ は対称行列全体)が $\dim M_{2,2}$ に
一致することを確かめよ。
\end{mondai}

\begin{mondai}[標準]\label{q:dim-ex8}
$\vect{a}_1 = \mat{1 \\ 2 \\ 3}$, $\vect{a}_2 = \mat{2 \\ 3 \\ 4}$,
$\vect{a}_3 = \mat{3 \\ 5 \\ 7}$ は $\R^3$ の基底か。基底でない場合は、
$W = \langle \vect{a}_1, \vect{a}_2, \vect{a}_3 \rangle$ の基底を
この中から選び出し、$\dim W$ を答えよ。
\end{mondai}

\begin{mondai}[発展]\label{q:dim-ex9}
漸化式 $a_{n+2} = 3a_{n+1} - 2a_n$ を満たす数列全体を $W$ とする
($W$ がベクトル空間になることは問 \ref{q:vs-ex6} と同様に示せる)。\\
　(1) 定数数列 $(1, 1, 1, \dots)$ と等比数列 $(1, 2, 4, 8, \dots)$
($a_n = 2^n$)がどちらも $W$ に属することを確かめよ。\\
　(2) この $2$ 本が $W$ の基底であることを示し、初項が $a_0$, $a_1$ の
一般の数列 $(a_n) \in W$ の一般項を求めよ。
\end{mondai}

\begin{mondai}[発展]\label{q:dim-ex10}
$\R^3$ の部分空間
$W_1 = \{\,(x, y, z) \mid x + y + z = 0\,\}$,
$W_2 = \{\,(x, y, z) \mid x = y\,\}$ について、
$\dim W_1$, $\dim W_2$, $\dim (W_1 \cap W_2)$, $\dim (W_1 + W_2)$ を
それぞれ求め、
\[
\dim(W_1 + W_2) = \dim W_1 + \dim W_2 - \dim(W_1 \cap W_2)
\]
が成り立っていることを確かめよ(この等式は\textbf{次元公式}と呼ばれ、
一般に成り立つことが知られている)。
\end{mondai}

\begin{mondai}[証明]\label{q:dim-ex11}
$\vect{v}_1, \vect{v}_2, \vect{v}_3$ をベクトル空間 $V$ の基底とする。
このとき
\[
\vect{w}_1 = \vect{v}_1, \qquad
\vect{w}_2 = \vect{v}_1 + \vect{v}_2, \qquad
\vect{w}_3 = \vect{v}_1 + \vect{v}_2 + \vect{v}_3
\]
も $V$ の基底であることを証明せよ。
\end{mondai}

\begin{mondai}[証明]\label{q:dim-ex12}
有限次元ベクトル空間 $V$ の部分空間 $W_1, W_2$ について
\[
\dim(W_1 + W_2) \le \dim W_1 + \dim W_2
\]
を証明せよ(ヒント:$W_1$ の基底と $W_2$ の基底を合わせた組が
$W_1 + W_2$ を生成することを示し、間引き定理 \ref{thm:dim-extend} を使う)。
\end{mondai}

\bigskip
{\small
\noindent\textbf{【章末問題の方針】}\quad
まず自力で取り組み、行き詰まったら下の方針を見よ。記述による完全解答は
巻末「問の解答」にある。
\begin{itemize}
\item \textbf{問 \ref{q:dim-ex1}}:$y$ について解き、残り $2$ 変数を
パラメータにする(例題 \ref{rei:dim-plane} と同じ手順)。
\item \textbf{問 \ref{q:dim-ex2}}:行列式で独立を確かめ、
定理 \ref{thm:dim-ncheck} を使う。座標は連立方程式で。
\item \textbf{問 \ref{q:dim-ex3}}:対角行列を「基本の対角行列」の
線形結合で書いてみる。
\item \textbf{問 \ref{q:dim-ex4}}:$W$ の元は $a_1 x + a_2 x^2$ の形
(問 \ref{q:sub-ex4})。
\item \textbf{問 \ref{q:dim-ex5}}:$x_1$ について解き、$x_2, x_3, x_4$ を
パラメータにする。
\item \textbf{問 \ref{q:dim-ex6}}:独立の確認は係数比較。
$\dim \R[x]_2 = 3$ と $n$ 本判定で基底と分かる。座標も係数比較で。
\item \textbf{問 \ref{q:dim-ex7}}:${}^t\!A = -A$ から対角成分が $0$ に
なることを導く。残る自由度を数える。
\item \textbf{問 \ref{q:dim-ex8}}:行列式を計算する。$0$ なら従属。
どれか $1$ 本を残り $2$ 本で表してみる。
\item \textbf{問 \ref{q:dim-ex9}}:例題 \ref{rei:dim-seq} と同じ筋書き。
「最初の $2$ 項が一致すれば全体が一致する」が鍵。
\item \textbf{問 \ref{q:dim-ex10}}:$W_1 \cap W_2$ は連立して解く。
$W_1 + W_2$ は $W_1$ より真に大きいことを示し、
定理 \ref{thm:dim-sub} で次元を確定する。
\item \textbf{問 \ref{q:dim-ex11}}:独立だけ示せば $n$ 本判定が使える。
$c_1\vect{w}_1 + c_2\vect{w}_2 + c_3\vect{w}_3 = \vect{0}$ を
$\vect{v}_i$ でまとめ直す。
\item \textbf{問 \ref{q:dim-ex12}}:和空間の任意の元
$\vect{w}_1 + \vect{w}_2$ を、二組の基底の線形結合で書く。
\end{itemize}
}

% =====================================================
\chapter{線形写像 ― 構造を保つ写像}\label{chap:linmap}
% =====================================================

\section{定義と例}

\begin{teigi}[線形写像]
ベクトル空間 $V$ から $W$ への写像 $f$ が、任意の
$\vect{u}, \vect{v} \in V$、$c \in \R$ に対して
\[
f(\vect{u} + \vect{v}) = f(\vect{u}) + f(\vect{v}),
\qquad
f(c\vect{u}) = c\,f(\vect{u})
\]
を満たすとき、$f$ を\textbf{線形写像}という。
\end{teigi}

第I部で「行列 $A$ の定める写像 $\vect{x} \mapsto A\vect{x}$」として
出会った概念の、抽象版である。新しい舞台では、驚くべき例が加わる。

\begin{rei}[微分は線形写像である]
微分 $D : \R[x]_n \to \R[x]_n$、$D(p) = p'$ は
\[
(p + q)' = p' + q', \qquad (cp)' = c\,p'
\]
を満たすから線形写像である。高校で「項別に微分してよい」と教わった規則は、
「微分は線形写像である」という一文に凝縮される。
定積分 $p \mapsto \int_0^1 p(x)\,dx$($\R[x]_n \to \R$)も線形写像である。
\end{rei}

\begin{rei}[数列のシフト]
数列空間上で、$(a_0, a_1, a_2, \dots) \mapsto (a_1, a_2, a_3, \dots)$ と
一つずらす写像も線形写像である。漸化式の理論はこの写像の分析にほかならない。
\end{rei}

\section{核と像、そして次元定理}

\begin{teigi}[核と像]
線形写像 $f : V \to W$ に対して、
\[
\mathrm{Ker}\,f = \{\, \vect{v} \in V \mid f(\vect{v}) = \vect{0} \,\},
\qquad
\mathrm{Im}\,f = \{\, f(\vect{v}) \mid \vect{v} \in V \,\}
\]
をそれぞれ $f$ の\textbf{核}、\textbf{像}という。
$\mathrm{Ker}\,f$ は $V$ の、$\mathrm{Im}\,f$ は $W$ の部分空間になる。
\end{teigi}

$f(\vect{x}) = A\vect{x}$ のとき、$\mathrm{Ker}\,f$ は同次方程式の解空間であり、
$\mathrm{Im}\,f$ は $A$ の列ベクトルが生成する空間(その次元は $\mathrm{rank}\,A$)である。
つまり核と像は、第II部の主役たちの抽象化である。

\begin{teiri}[単射性の判定]
線形写像 $f$ が単射(一対一)であることと、
$\mathrm{Ker}\,f = \{\vect{0}\}$ は同値である。
\end{teiri}

「写像がつぶれていないこと」を、たった一点 $\vect{0}$ の逆像だけで
判定できる――線形性のご利益である
($f(\vect{u}) = f(\vect{v}) \iff f(\vect{u} - \vect{v}) = \vect{0}$ に注意)。

\begin{teiri}[次元定理]
$V$ が有限次元のとき、線形写像 $f : V \to W$ に対して
\[
\dim V = \dim \mathrm{Ker}\,f + \dim \mathrm{Im}\,f.
\]
\end{teiri}

$f(\vect{x}) = A\vect{x}$($A$ は $m \times n$)の場合に翻訳すると
\[
n = (n - \mathrm{rank}\,A) + \mathrm{rank}\,A
\]
となり、第II部の自由度定理が再現される。次元定理は
「入力の次元 = つぶれて消える次元 + 生き残って出力される次元」という
\textbf{情報の保存則}として読める。

\section{同型 ― 「見た目が違うだけ」を厳密にいう}

\begin{teigi}[同型]
線形写像 $f : V \to W$ が全単射であるとき $f$ を\textbf{同型写像}といい、
同型写像が存在する $V$ と $W$ は\textbf{同型}であるという($V \cong W$ と書く)。
\end{teigi}

\begin{teiri}
$n$ 次元ベクトル空間 $V$ は、すべて $\R^n$ と同型である。
実際、基底を一組固定し、各ベクトルをその座標にうつす写像
$\vect{v} \mapsto (c_1, \dots, c_n)$ が同型写像になる。
\end{teiri}

\begin{itembox}[l]{深掘り:同型の哲学 ― では抽象論は無駄だったのか}
「どうせすべて $\R^n$ と同型なら、最初から $\R^n$ だけ考えれば
よかったのでは?」――もっともな疑問だが、答えは否である。
理由は二つある。第一に、$\R^n$ との同型を作るには\textbf{基底を選ぶ}必要があり、
その選び方は一通りではない。多項式の空間に「自然な」基底 $1, x, x^2, \dots$ が
あるのは偶然で、一般の空間(たとえば微分方程式の解空間)には
特権的な基底は存在しない。基底に依存しない概念(次元、核、像、
後に学ぶ固有値)と、基底に依存する表示(座標、成分)を区別する目が、
抽象論によって養われる。第二に、無限次元空間($C(\R)$ など)は
そもそも $\R^n$ に帰着できない。量子力学や関数解析はその世界の話であり、
有限次元の抽象論はそこへの登山口である。
\end{itembox}

\begin{mondai}\label{q:linmap}
微分 $D : \R[x]_3 \to \R[x]_3$ について、$\mathrm{Ker}\,D$ と
$\mathrm{Im}\,D$ を求め、それぞれの次元が次元定理を満たすことを確かめよ。
\end{mondai}

% =====================================================
\chapter{表現行列と基底の取り替え}\label{chap:repmat}
% =====================================================

\section{線形写像を行列で書く}

抽象的な線形写像も、基底を選んで座標で見れば、結局は行列になる。

\begin{teigi}[表現行列]
$V$ の基底 $\vect{e}_1, \dots, \vect{e}_n$ をとる。
線形写像 $f : V \to V$ に対し、各基底ベクトルの行き先を基底で展開して
\[
f(\vect{e}_j) = a_{1j}\vect{e}_1 + a_{2j}\vect{e}_2 + \cdots + a_{nj}\vect{e}_n
\]
と書いたときの係数を並べた行列 $A = (a_{ij})$ を、
この基底に関する $f$ の\textbf{表現行列}という。
\textbf{第 $j$ 列には $f(\vect{e}_j)$ の座標が入る}と覚える。
\end{teigi}

座標の言葉では、$\vect{v}$ の座標を $\vect{c}$ とすると
$f(\vect{v})$ の座標は $A\vect{c}$ になる。
抽象的な写像の計算が、行列の掛け算に完全に翻訳される。

\begin{rei}[微分の表現行列]
$D : \R[x]_2 \to \R[x]_2$ を基底 $1, x, x^2$ で表す。
\[
D(1) = 0, \qquad D(x) = 1, \qquad D(x^2) = 2x
\]
の座標を列に並べて
\[
A = \mat{0 & 1 & 0 \\ 0 & 0 & 2 \\ 0 & 0 & 0}.
\]
たとえば $p(x) = 3 + 5x + 4x^2$(座標 $(3, 5, 4)$)に対し
\[
A \mat{3 \\ 5 \\ 4} = \mat{5 \\ 8 \\ 0}
\]
は $p'(x) = 5 + 8x$ の座標である。\textbf{微分という解析の操作が、
行列の掛け算という代数の操作になった}。
なお $A^3 = O$ となる($3$ 回微分すると $2$ 次以下の多項式は消える)。
$A \neq O$ なのに $A^3 = O$ となる行列(\textbf{べき零行列})は、
第VI部のジョルダン標準形で主役を演じる。
\end{rei}

\section{基底を取り替えると行列はどう変わるか}

同じ線形写像でも、基底を取り替えれば表現行列は変わる。
新基底 $\vect{e}'_1, \dots, \vect{e}'_n$ を旧基底で展開した係数を
列に並べた行列 $P$(\textbf{基底変換行列}、必ず正則)を使うと、
新旧の座標は $\vect{c} = P\vect{c}'$ で結ばれ、表現行列は

\begin{teiri}[表現行列の変換則]
旧基底での表現行列を $A$、新基底での表現行列を $A'$ とすると
\[
A' = P^{-1} A P.
\]
\end{teiri}

$A' = P^{-1}AP$ の関係にある二つの行列は\textbf{相似}であるという。
相似な行列は「同じ変換を別の座標系から見たもの」にすぎない。

\begin{itembox}[l]{深掘り:座標は人間の都合、変換は実在 ― そして不変量}
回転や微分という\textbf{変換そのもの}は、人間が座標をどう引くかに
関係なく存在する。座標と表現行列は、それを記述するための
\textbf{観測者の都合}である。だとすれば、相似で変わらない量
($= $ どの座標系から見ても同じ量)こそが、変換の本当の性質を表すはずである。
実際、
\[
\det(P^{-1}AP) = \det P^{-1} \det A \det P = \det A
\]
より\textbf{行列式は相似不変}である(拡大率が座標の取り方によらないのは当然だ)。
対角成分の和 $\mathrm{tr}\,A = a_{11} + \cdots + a_{nn}$(\textbf{トレース})も
相似不変であることが知られている。階数も、そして第IV部で学ぶ
\textbf{固有値}も相似不変である。

逆に問おう。座標の取り方は自由なのだから、
\textbf{表現行列が最も簡単になる基底}を選べばよいのではないか?
最良の場合、表現行列は対角行列になる――これが\textbf{対角化}であり、
第IV部の主題である。第III部の抽象論は、この問いを立てるための準備だった。
\end{itembox}

\begin{mondai}\label{q:repmat}
$\R^2$ の線形変換 $f$ を標準基底に関する表現行列
$A = \mat{3 & 1 \\ 1 & 3}$ で定める。
新基底 $\vect{e}'_1 = \mat{1 \\ 1}$, $\vect{e}'_2 = \mat{1 \\ -1}$
に関する表現行列 $A' = P^{-1}AP$ を計算せよ
($P$ は新基底を列に並べた行列)。結果が対角行列になることを確かめよ。
\end{mondai}
